% Refactored five-act IEEE Transactions main text for Performance Before Outcomes.

\section{Introduction}\label{sec:intro}

\IEEEPARstart{H}{ominin} endocranial volume roughly tripled over the past three million years, rising from Australopithecine means near $450\ccm$ to late Pleistocene means near $1{,}350$--$1{,}500\ccm$ (Fig.~\ref{fig:main}A)~\cite{du2018,will2021,gingerich2022}. Across $\sim\!10^5$ generations, that net trajectory requires a long-run directional response of only $\sim\!10^{-4}$ phenotypic standard deviations ($\SD$) per generation---two to three orders of magnitude smaller than selection differentials routinely measured in living populations~\cite{kingsolver2001}. Yet brain tissue is metabolically expensive, consuming roughly a fifth of resting adult energy and peaking at \PboChildPeakRmr{} of resting metabolism in early childhood~\cite{pontzer2016,kuzawa2014}. Meanwhile, human children require nearly two decades of biparental and alloparental provisioning~\cite{kaplan2000,marlowe2000}, and coalitionary cooperation unfolds across years of reciprocal exchange~\cite{boehm1999,gurven2004}. Prospective partners and band members must therefore commit \emph{before} completed family size, crisis provisioning, or lifelong fidelity can be observed---selecting on \emph{performance before outcomes}.

This structural timing problem links two major unresolved gaps across evolutionary anthropology and modern partner choice:
\begin{enumerate}[leftmargin=1.4em,itemsep=1.5pt,topsep=2pt]
\item \textbf{When the human brain grew, when it stalled, and what released it (Section~\ref{sec:fossils}):} Why did endocranial volume expand between $2.1$ and $1.5\,\mathrm{Ma}$, stall near $\sim\!\PboPlateauMeanCC\ccm$ for $800\,\mathrm{kyr}$ ($1.50$--$0.70\,\mathrm{Ma}$) hundreds of millennia \emph{before} Mid-Pleistocene climate crises or proposed bottlenecks, remain small in insular or calorie-constrained \Homo{} lineages, and then surge $+50\%/\mathrm{Myr}$ simultaneously across Africa and Eurasia after $0.7\,\mathrm{Ma}$?
\item \textbf{Secondary entry filters versus primary pair-bond criteria (Sections~\ref{sec:juries}--\ref{sec:platforms}):} Why do single-context stranger markets select heavily on \emph{secondary mating criteria}---attributes with high zero-acquaintance legibility ($\rho_X$) such as static portrait rank, $6'0''$ stature cutoffs, income brackets, hobby tags, and demographic sorting labels that contribute near zero to union satisfaction or completed fertility ($\alpha_X\approx 0$)---whereas \emph{primary mating criteria} ($\rho_{K,0}\ll\rho_X$: everyday kindness, conversational play, general intelligence $g$, mutual reverence, emotional fidelity, and unstaged hearth devotion) require multi-context witness breadth yet govern pair-bond love, survival, and juvenile encephalization ($\alpha_K,\alpha_D\gg|\alpha_X|$)?
\end{enumerate}

Appendix~\ref{app:theory} formalizes this architecture in twenty-four propositions across three stages: (i)~\emph{what observation reveals} across domains and contexts (Propositions~\ref{prop:identity}--\ref{prop:division},~\ref{prop:twoway}), proving that $K$-domain residential breadth multiplies general-capacity aim $K$-fold, taxes single-domain gaming as $1/K$, and breaks the single-context validity ceiling $R^2_{\infty,1}<1$; (ii)~\emph{what sustains delivery, love, marital survival, and legitimate standing} (Propositions~\ref{prop:despot},~\ref{prop:delivery}--\ref{prop:enforce},~\ref{prop:love},~\ref{prop:lifecycle},~\ref{prop:prestige}), proving that staged tests ($d\to 1$) erase delivery validity $R_p\to 0$, stranger legibility inverts secondary entry weights against primary pair-bond welfare, compressed pre-cohabitation witness audits inflate lock-in failure and asymmetric marital exit, and multiplicative standing ($\widehat W=\hat b\hat p$) separates earned outlier confidence from gross arrogance across audience priors and executive sex asymmetries (Section~\ref{sec:prestige}); and (iii)~\emph{how juvenile metabolism and pairing gates allocate evolution} (Propositions~\ref{prop:persist}--\ref{prop:episodes},~\ref{prop:ancestors}--\ref{prop:waiting},~\ref{prop:childhood}), proving that cooking efficiency and pair-bond devotion break the solitary-mother childhood brain ceiling supermodularly while modern reproduction is gated by union entry rather than marital birth ideals.

Rather than estimating unobservable Pleistocene selection gradients, we test the observable comparative statics of these equations across fifteen open and archival datasets spanning deep time, cross-cultural ethnography and song corpora, laboratory encounters, longitudinal and national couples cohorts, deep-learning facial attribution, and \PboThreePlatformRecords{} digital platform records (including \PboOkcEssaysN{} multi-domain \texttt{OkCupid} essays and \PboVoiceEssaysN{} Plenty of Fish self-descriptions). Table~\ref{tab:macro_kpis} summarizes the nine macro selection key performance indicators (KPIs) and scope boundaries.

\begin{table*}[!t]
\caption{Executive Scorecard: Nine Macro Selection KPIs Across Brain Evolution, Secondary vs.\ Primary Mating Criteria, Prestige, and Modern Platforms.}\label{tab:macro_kpis}
\scriptsize
\renewcommand{\arraystretch}{0.96}
\setlength{\tabcolsep}{3.5pt}
\begin{tabularx}{\textwidth}{@{}>{\raggedright\arraybackslash}p{0.165\textwidth}>{\raggedright\arraybackslash}p{0.200\textwidth}>{\raggedright\arraybackslash}p{0.145\textwidth}LL@{}}
\toprule
Macro Selection KPI & Core Selection Question & Empirical Scale & Headline Quantitative Contrast & Boundary Check \& Formal Link \\
\midrule
\tabinput{performance-macro-kpi-table.tex}
\bottomrule
\end{tabularx}
\tablenote{Every cell is generated programmatically from SHA-256-verified outputs (\texttt{scripts/run.py --study performance-before-outcomes}). Observational and contested tiers are labeled explicitly; pre-committed falsifiers F1--F16 appear in Appendix Table~\ref{tab:falsifiers}.}
\end{table*}

\begin{figure*}[!t]
\centering
\includegraphics[width=\textwidth]{performance-fig-fossils.pdf}
\caption{\textbf{When the hominin brain grew, why it stalled for $800\,\mathrm{kyr}$, and how cooking, childhood pair-bond care, and hearth song legibility released it.} \textbf{(A)}~Endocranial volume across \FosN{} crania~\cite{gingerich2022,du2018,will2021} ($\Delta\mathrm{AICc}=0$ stasis--step vs.\ $\PboLinearDeltaAICc$ log-linear) synchronized with the 57-core \texttt{LR04} $\delta^{18}\mathrm{O}$ stack~\cite{lisiecki2005}: $1.50$--$0.70\,\mathrm{Ma}$ stasis ($\sim\!\PboPlateauMeanCC\ccm$) preceded $1.12\,\mathrm{Ma}$ cooling~\cite{margari2023}, the contested $930$--$813\,\mathrm{ka}$ bottleneck~\cite{hu2023,deng2025}, and $789\,\mathrm{ka}$ tektites~\cite{jourdan2019} by $\PboGapCoolingKyr$--$\PboGapTektiteKyr\,\mathrm{kyr}$ before $1.0$--$0.78\,\mathrm{Ma}$ fire and cooking~\cite{berna2012,zohar2022}. \textbf{(B)}~Epistemic audit of ten $1.25$--$0.70\,\mathrm{Ma}$ events and within-clade slowdown ($\PboDuRatePre\to\PboDuRatePlateau\,\ln/\mathrm{Myr}$, $-\PboDuRateDropPct$~\cite{du2018}). \textbf{(C)}~Age-4.5 brain metabolic peak (\PboChildPeakRmr{} RMR, $\PboChildAdultRmrRatio\times$ adult~\cite{kuzawa2014}), \PboSccsChildSoloMomPct{} solo-mother care across 136 SCCS societies~\cite{isler2012,hrdy2009}, and Proposition~\ref{prop:childhood} Lande equilibria ($\PboETwentyTwoSolo\to\PboETwentyTwoRawAllo\to\PboETwentyTwoCookNoLove\to\PboETwentyTwoFull\ccm$; $\PboETwentyTwoLoveGain\ccm$ pair-bond gain). \textbf{(D,~E)}~Post-$0.7\,\mathrm{Ma}$ $+50\%/\mathrm{Myr}$ multi-regional surge~\cite{will2021} alongside small-brained \textit{H.~naledi}~\cite{berger2015} and \textit{H.~floresiensis}~\cite{kubo2013}. \textbf{(F)}~Harvard \texttt{Natural History of Song} (\PboNhsSongsN{} recordings, \PboNhsSocN{} societies, \PboNhsRatersN{} blind listeners~\cite{mehr2019}): lullabies ($d=\PboNhsLullabyD$) and communal dance ($d=\PboNhsDanceD$) are $\PboNhsLullabyLoveRatio\times$ and $\PboNhsDanceLoveRatio\times$ more stranger-legible than love songs ($d=\PboNhsLoveD$).}\label{fig:fossils}
\end{figure*}

\section{When the Human Brain Grew, Why It Stalled for \texorpdfstring{800\,kyr}{800 kyr}, and What Released It}\label{sec:fossils}

\subsection{Four Regimes of Growth and Stasis---and the Onset--Crisis Timing Constraint}
Across the three open hominin cranial compilations---\FosN{} specimens over $3.19\,\mathrm{Myr}$ in Gingerich~\cite{gingerich2022}, 166 Pleistocene \Homo{} crania in Will et al.~\cite{will2021}, and 94 specimens in Du et al.~\cite{du2018}---a constant log-linear expansion is decisively rejected ($\Delta\mathrm{AICc}=\PboLinearDeltaAICc$; Fig.~\ref{fig:fossils}A). Instead, hominin brain evolution alternated sharply between stasis ceilings and rapid releases across four regimes: (i)~an Australopithecine ceiling near $\FosLevI\ccm$ prior to $\sim\!2.1\,\mathrm{Ma}$ (Phase I, when brains did not grow despite bipedality); (ii)~an early-\Homo/\textit{H.~erectus} expansion doubling volume to $\sim\!\PboPlateauMeanCC\ccm$ between $2.1$ and $\PboPlateauStartFit\,\mathrm{Ma}$ (Phase II); (iii)~an $\sim\!800\,\mathrm{kyr}$ Middle/Early Pleistocene stasis plateau between $\sim\!1.50$ and $0.70\,\mathrm{Ma}$ (Phase III, within-window slope $+0.02\pm0.11\,\ln/\mathrm{Myr}$, indistinguishable from zero; within-clade Monte Carlo rate dropping $-\PboDuRateDropPct$ from $\PboDuRatePre$ to $\PboDuRatePlateau\,\ln/\mathrm{Myr}$~\cite{du2018}); and (iv)~a renewed $+50\%/\mathrm{Myr}$ multi-continental surge after $\sim\!\PboPlateauEndFit\,\mathrm{Ma}$ (Phase IV), even as calorie- or island-constrained \textit{H.~naledi} ($465$--$560\ccm$) and \textit{H.~floresiensis} ($426\ccm$) retained small brains into the Late Pleistocene.

Synchronizing this cranial record with the 57-core \texttt{LR04} benthic foraminiferal $\delta^{18}\mathrm{O}$ stack~\cite{lisiecki2005} establishes a decisive chronological constraint on catastrophe-based theories of why the brain stalled (Fig.~\ref{fig:fossils}A,B; Appendix Table~\ref{tab:plateau_events}). Phase III stasis began near $1.50\,\mathrm{Ma}$ (fitted $\PboPlateauStartFit\,\mathrm{Ma}$), under mild 41\,kyr obliquity cycles (\texttt{LR04} P90--P10 swing $\PboMptSwingPre\text{\textperthousand}$). Every major environmental catastrophe invoked in recent debates occurred hundreds of millennia \emph{after} stasis was already established: $+\PboGapMptKyr\,\mathrm{kyr}$ for Mid-Pleistocene Transition onset ($1.25\,\mathrm{Ma}$), $+\PboGapCoolingKyr\,\mathrm{kyr}$ for MIS~34 European cooling ($1.12\,\mathrm{Ma}$), $+\PboGapBottleKyr\,\mathrm{kyr}$ for the contested $930$--$813\,\mathrm{ka}$ coalescent bottleneck, and $+\PboGapTektiteKyr\,\mathrm{kyr}$ for the $789\,\mathrm{ka}$ Australasian impact. Under Proposition~\ref{prop:tissue} and Corollary~\ref{cor:colimit}, early stasis onset without catastrophe reflects energetic and juvenile co-limitation: once raw-foraging gut allocation and \textit{H.~erectus} body size hit pre-cooking caloric limits ($\kappa_B\bar B\approx a\beta_{\mathrm{eff}}$)~\cite{aiello1995,navarrete2011,pontzer2016}, brain volume stalled near $\PboPlateauMeanCC\ccm$ despite ongoing social selection.

\subsection{Separating Late-Plateau Crises from the Cooking, Childhood, and Hearth Release}
Evaluating overlapping 1.25--0.70\,Ma global events by epistemic tier separates processes that constrained late-plateau habitats or triggered the 0.7\,Ma release from contested or unlinked hypotheses (Fig.~\ref{fig:fossils}B; Table~\ref{tab:plateau_events}):
\begin{itemize}[leftmargin=1.3em,itemsep=1.5pt,topsep=2pt]
\item \textbf{Measured paleoclimatic cooling and regional depopulation:} Across the Mid-Pleistocene Transition (1.25--0.70\,Ma)~\cite{clark2006}, glacial periodicity shifted from 41\,kyr to high-amplitude 100\,kyr cycles (\texttt{LR04} P90--P10 glacial--interglacial swings rising $51\%$ from $\PboMptSwingPre\text{\textperthousand}$ to $\PboMptSwingPost\text{\textperthousand}$), accompanied by atmospheric $\mathrm{CO}_2$ drawdown~\cite{chalk2017,yan2019} and a severe 950--860\,ka (MIS 24--22) ocean-circulation crisis~\cite{pena2014}. Earlier, at MIS 34 ($\sim\!1.12\,\mathrm{Ma}$, \texttt{LR04} peak $\PboMisThirtyFourD\text{\textperthousand}$), climate-envelope modeling and Iberian marine pollen cores indicate extreme cooling that plausibly drove regional European hominin depopulation~\cite{margari2023} prior to Po Valley/southern European recolonization near 0.9\,Ma~\cite{muttoni2024}. Crucially, this depopulation was \emph{regional} to western Eurasia rather than a global extinction, while climate-driven habitat fragmentation later accelerated allopatric diversification~\cite{timmermann2022,ao2026}.
\item \textbf{Contested coalescent bottleneck (930--813\,ka):} Using \texttt{FitCoal} on modern site-frequency spectra, Hu et al.~\cite{hu2023} and Zhou et al.~\cite{zhou2026} inferred a prolonged ancestral contraction to coalescent effective size $N_e\approx 1{,}280$, proposing links to chromosome~2 fusion and Phase IV speciation. Two qualifications are essential. First, $N_e\approx 1{,}280$ is a harmonic-mean coalescent parameter shaped by population structure and selective sweeps, not a literal headcount of living individuals (Proposition~\ref{prop:ancestors}). Second, independent statistical reanalyses found that alternative demographic histories without a severe bottleneck fit modern SFS and linkage patterns as well or better, attributing the \texttt{FitCoal} signal to numerical smoothing artifacts or weakly structured African stem populations~\cite{deng2025,cousins2025,terhorst2025,ragsdale2023,cousins2024}, consistent with continuous 930--813\,ka Acheulean occupation across Africa, the Levant, and Spain~\cite{goreninbar2000,pares2013}. We therefore mark the severe bottleneck as a \emph{contested genomic hypothesis}.
\item \textbf{Impacts, super-eruptions, and magnetic reversal without hominin die-off:} Five dated geophysical events overlap Phase III---the $\sim\!1.07\,\mathrm{Ma}$ Bosumtwi impact (10.5\,km crater, Ghana)~\cite{koeberl2007}, $\sim\!815\,\mathrm{ka}$ Pantasma impact (14\,km, Nicaragua)~\cite{rochette2019}, $\sim\!800\,\mathrm{ka}$ Oldest Toba Tuff super-eruption~\cite{ito2020}, $\sim\!789\,\mathrm{ka}$ Australasian tektite impact scattered across 10--30\% of Earth's surface (linked to a buried Bolaven Plateau source in Laos~\cite{jourdan2019,sieh2020,sieh2024}, though crater identification remains debated~\cite{mizera2024}), and the $\sim\!773\,\mathrm{ka}$ Matuyama--Brunhes geomagnetic reversal~\cite{singer2019}. None establishes a global hominin die-off; notably, 789\,ka Bose Basin biface assemblages in southern China and Vietnam show behavioral continuity directly across the tektite horizon~\cite{michel2024}.
\item \textbf{Campfire energetic release, childhood pair-bond ceiling, and hearth song legibility:} Controlled fire at Wonderwerk Cave ($\sim\!1.0\,\mathrm{Ma}$)~\cite{berna2012} and structured fish cooking at Gesher Benot Ya'aqov ($\sim\!780\,\mathrm{ka}$)~\cite{zohar2022,wrangham2009} align directly with the $\sim\!\PboPlateauEndFit\,\mathrm{Ma}$ termination of Phase III. Crucially, the sharpest metabolic bottleneck sits in childhood rather than adulthood: PET/MRI glucose data show that the human brain peaks at \PboChildPeakRmr{} of resting metabolic rate (\PboChildPeakDer{} of daily energy expenditure) at age $4.5$---$\PboChildAdultRmrRatio\times$ the adult share ($20.5\%$)~\cite{kuzawa2014}, imposing the $\sim\!650\ccm$ primate ``Gray Ceiling'' under solitary maternal foraging~\cite{isler2012,hrdy2009} (Fig.~\ref{fig:fossils}C). Across 136 coded Standard Cross-Cultural Sample societies~\cite{kirby2016}, mothers provide exclusive early-childhood care in \PboSccsChildSoloMomPct{} of societies ($0/136$; non-mothers provide majority care in \PboSccsChildAlloMajPct{}), paternal closeness rises from \PboSccsFatherInfPct{} in infancy to \PboSccsFatherChildPct{} in early childhood (\PboSccsFatherOccPct{} at least occasional), and community food sharing operates in \PboSccsFoodSharePct{}. Under Proposition~\ref{prop:childhood} (\texttt{E22}, verified across 20,000 simulated agents within $\le 0.75\ccm$), Lande equilibria step from $\PboETwentyTwoSolo\ccm$ (solitary ape mother) to $\PboETwentyTwoRawAllo\ccm$ (raw alloparental plateau), stall at $\PboETwentyTwoCookNoLove\ccm$ under cooking \emph{without} durable pair-bond commitment, and reach $\PboETwentyTwoFull\ccm$ only when cooking combines supermodularly with pair-bond devotion ($\PboETwentyTwoLoveGain\ccm$ love gain). Evening hearths simultaneously extended multi-domain social auditing ($K\uparrow$, $(1-d)n\uparrow$): across \PboNhsSongsN{} field recordings from \PboNhsSocN{} societies rated blind by \PboNhsRatersN{} naive listeners in the Harvard \texttt{Natural History of Song} corpus~\cite{mehr2019} (\PboNhsEthPerfN{} ethnographic performances), zero-context strangers decode infant lullabies ($d=\PboNhsLullabyD$, AUROC $0.89$) and communal motor dance ($d=\PboNhsDanceD$, AUROC $0.84$; \PboNhsDanceShare{} of coded hearth performances) $\PboNhsLullabyLoveRatio\times$ and $\PboNhsDanceLoveRatio\times$ more strongly than romantic love songs ($d=\PboNhsLoveD$, AUROC $0.61$; Fig.~\ref{fig:fossils}F)---confirming that alloparental soothing and coalitionary synchrony are species-legible public broadcasts whereas romantic devotion requires lived relational proof. Immediately post-0.7\,Ma, African \Homo{} into \Hs{} ($\SerAfr\,\ln/\mathrm{Myr}$)~\cite{hublin2017}, western Eurasian \Homo{} into Neanderthals ($\SerEur\,\ln/\mathrm{Myr}$), and East Asian archaic \Homo{} ($\SerAsia\,\ln/\mathrm{Myr}$)~\cite{ni2021} expanded in parallel at $+50\%/\mathrm{Myr}$ (Fig.~\ref{fig:fossils}D), whereas small-brained \textit{H.~naledi} ($465$--$560\ccm$ at $\sim\!285\,\mathrm{ka}$)~\cite{berger2015} and insular \textit{H.~floresiensis} ($426\ccm$ at $\sim\!80\,\mathrm{ka}$)~\cite{kubo2013} persisted where local ecological budgets kept $\kappa_B$ high (Fig.~\ref{fig:fossils}E).
\end{itemize}

\begin{figure*}[!t]
\centering
\includegraphics[width=\textwidth]{performance-fig-jury-choice.pdf}
\caption{\textbf{Four-minute stranger legibility, the two-way jury law, unique dyadic spark, the 17-hobby checklist null, and longitudinal hindsight.} \textbf{(A)}~Stated pre-event weights (open) versus revealed fixed-effects choice weights (filled) across \SdRows{} speed dates of \SdN{} participants~\cite{fisman2006}. \textbf{(B)}~Single-observer agreement ($r_1$) governs revealed choice slope: looks ($r_1=\JurOneAtt$) are immediately legible, whereas intelligence ($r_1=\JurOneInt$) and sincerity ($r_1=\JurOneSin$) are nearly opaque to one stranger. \textbf{(C)}~Two-way validity $R^2_{m,k}$ across $m$ observers and $k$ contexts (Proposition~\ref{prop:twoway}): single-context crowds ($k=1$) stall at $R^2_{\infty,1}=\PboTwoWaySingleAsym$, whereas a residential band ($m=8,k=8$) reaches $R^2_{8,8}=\PboTwoWayEightEight$. \textbf{(D)}~Bipartite Social Relations Model ($n=\PboDyadLikeN$ ratings, $\PboDyadPairsN$ dyads~\cite{eastwick2014}): partner consensus explains only \PboDyadConsensusShare{} of liking (\PboDyadIntelConsensus{} intelligence, \PboDyadSincConsensus{} sincerity), whereas unique pair-specific spark accounts for \PboDyadUniqueShare{} ($r=\PboDyadReciprocityR$ reciprocity). \textbf{(E)}~Across \PboHobbyDyadsN{} dyads, objective 17-hobby tag overlap (\texttt{int\_corr}) is null for ``Yes'' ($r=\PboHobbyObjR$), whereas live conversational wavelength (\texttt{shar}) reaches $r=\PboHobbyLiveR$ ($\PboHobbyRatio\times$ stronger~\cite{finkel2012}). \textbf{(F)}~Waves 1--3 hindsight~\cite{fisman2006,eastwick2008}: 24\,h post-event, women's looks weight jumps $\PboHindLooksJumpW$ points before intelligence rebounds $\PboHindIntelSnapW$ points 3--4 weeks later.}\label{fig:jury}
\end{figure*}

\begin{figure*}[!t]
\centering
\includegraphics[width=\textwidth]{performance-fig-generalist-verify.pdf}
\caption{\textbf{Generalist breadth, anti-gaming dilution, staged versus unstaged delivery (\texttt{E15}--\texttt{E16}), and nine-domain mind selection.} \textbf{(A)}~Events per domain $n^*(K)$ to reach $R^*=0.5$ across $K=1\dots 24$ domains (Proposition~\ref{prop:time}): impossible for $K\le 5$, $\PolyNSix$ at $K=6$, and $\approx 1$ at $K=12$. \textbf{(B)}~Equilibrium general-capacity investment share $\pi_G=Kc_S/(Kc_S+c_G)$ (Proposition~\ref{prop:polymath}) and $1/K$ dilution of single-domain gaming (Corollary~\ref{cor:inflation}). \textbf{(C)}~Opposite response of capacity validity $R_b(d)$ and delivery validity $R_p(d)$ as staged share $d$ rises across $n=\VerN$ occasions (Proposition~\ref{prop:staging}). \textbf{(D)}~Equilibrium capacity and voluntary commitment across third-party enforced sharing $e\in[0,1]$ below ($T=\EnfT$) versus above ($T=\EnfCrowdT$) threshold (Proposition~\ref{prop:enforce}). \textbf{(E)}~Across \PboSatDatesN{} speed dates (\PboSatPartN{} participants~\cite{fisman2006}), one stranger explains \PboSatSingleRTwo{} of partner SAT ($r=\PboSatYesR$ for ``Yes'', $10+$ raters stall at \PboSatPanelRTwo{}), whereas observing \PboOkcDomWritersN{} \texttt{OkCupid-JSE} writers across $K=1\to 9$ prompts~\cite{kim2015} lifts lexical reliability $\PboOkcDomGain\times$ ($\PboOkcDomROne\to\PboOkcDomRNine$). \textbf{(F)}~Across writer lexical quintiles ($N=\PboOkcAssortN$~\cite{kim2015}), explicit requests for curiosity, intelligence, and wit climb $+40\%$ (\PboOkcAssortLo{}$\to$\PboOkcAssortHi{}) alongside kindness ($+38\%$), while height/money/looks tags stay flat ($\sim\!2.3\%$).}\label{fig:generalist_verify}
\end{figure*}

\section{Stranger Encounters: Mind-Reading Blindness, Multi-Context Juries, and Secondary Gaming}\label{sec:juries}

\subsection{Stranger Legibility, Cross-Sex Mind-Reading Blindness, and Post-Match Initiation}
Why do solitary first encounters overweight secondary surface cues while missing primary mind and devotion? Reanalyzing \SdRows{} four-minute Columbia speed-dating encounters among \SdN{} participants~\cite{fisman2006} exposes a sharp legibility wedge alongside systematic cross-sex misperception (Fig.~\ref{fig:jury}A,B; Fig.~\ref{fig:dynamics}A). Before meeting, women allocate nearly equal stated weights across intelligence, sincerity, fun, ambition, and looks (\PboLooksStatedW{} of 100 points to looks; men allocate \PboLooksStatedM). Yet when asked to predict what the \emph{opposite sex} seeks, both sexes severely overestimate opposite-sex preoccupation with secondary looks while underestimating primary mind and character: women overestimate men's looks weight by $\PboMindReadLooksW$ points ($95\%$ CI $[+5.9,+11.1]$) while underestimating men's valuation of intelligence ($\PboMindReadIntelW$ points, $[-7.9,-5.9]$) and sincerity ($\PboMindReadSincW$ points), and men mirror the error about women ($\PboMindReadLooksM$ points on looks, $\PboMindReadIntelM$ on intelligence, $\PboMindReadSincM$ on sincerity), reinforced by same-sex pluralistic ignorance ($\PboPeerIgnorLooksW$ points among women, $\PboPeerIgnorLooksM$ among men; Fig.~\ref{fig:dynamics}A). In actual four-minute ``Yes'' decisions with chooser fixed effects, both sexes act heavily on physical attractiveness (+$\PboLooksPanelW\,\mathrm{pp}/\SD$ in women, +$\PboLooksPanelM\,\mathrm{pp}/\SD$ in men; bottom-to-top quintile \PboYesLowW{}$\to$\PboYesHighW{} and \PboYesLowM{}$\to$\PboYesHighM{}).

Inter-observer reliability explains why without invoking hypocrisy (Fig.~\ref{fig:jury}B). Two strangers meeting a candidate for four minutes agree at $r_1=\JurOneAtt$ on looks and $r_1=\JurOneFun$ on fun, but only $r_1=\JurOneInt$ on intelligence and $r_1=\JurOneSin$ on sincerity---requiring $\sim\!6$ raters for looks versus $\sim\!25$ for intelligence and $>60$ for sincerity (Eq.~\eqref{eq:sb}). Proposition~\ref{prop:twoway} further proves that pooling $m\to\infty$ raters within a \emph{single} curated context ($k=1$) never averages out person-by-context state $\sigma_C^2$, stalling at $R^2_{\infty,1}=\PboTwoWaySingleAsym$ whereas a residential camp ($m=8,k=8$) reaches $R^2_{8,8}=\PboTwoWayEightEight$ (Fig.~\ref{fig:jury}C). Moreover, a finite-sample-unbiased bipartite Social Relations Model across $n=\PboDyadLikeN$ ratings and $\PboDyadPairsN$ dyads~\cite{eastwick2014} proves that attraction is relational rather than scalar ``trophy value'' (Fig.~\ref{fig:jury}D; Proposition~\ref{prop:love}(ii)): partner consensus explains only \PboDyadConsensusShare{} of liking (\PboDyadLooksConsensus{} looks, \PboDyadSharConsensus{} shared interests, \PboDyadIntelConsensus{} intelligence, \PboDyadSincConsensus{} sincerity), whereas chooser leniency explains \PboDyadChooserShare{} and unique pair-specific spark accounts for \PboDyadUniqueShare{} ($r=\PboDyadReciprocityR$ mutual residual reciprocity), lifting women's combined live spark slope ($\PboDyadSparkSlopeW\,\mathrm{pp}/\SD$; feeling chosen back $\PboDyadChosenBackW\,\mathrm{pp}$ W, $\PboDyadChosenBackM\,\mathrm{pp}$ M) above physical attractiveness ($\PboDyadLooksSlopeW\,\mathrm{pp}/\SD$).

Speed dating likewise refutes profile-tag matching~\cite{finkel2012,eastwick2008} (Fig.~\ref{fig:jury}E,F). Across \PboHobbyDyadsN{} dyads, exact objective correlation across 17 pre-scored leisure hobbies (\texttt{int\_corr}) has a null effect on saying ``Yes'' ($r=\PboHobbyObjR$, $[-0.019,+0.043]$; women's FE slope $\PboHobbyFeSlopeW\,\mathrm{pp}/\SD$), whereas live four-minute conversational wavelength (\texttt{shar}) correlates at $r=\PboHobbyLiveR$ ($\PboHobbyRatio\times$ stronger; $\PboHobbyLiveSlopeW\,\mathrm{pp}/\SD$). Tracking Waves 1--3 reveals both an initiation bottleneck and a primary-trait rebound: across $690$ mutual matches 3--4 weeks post-event, men place $\PboCallRatioMF\times$ as many follow-up calls per match as women ($49.7\%$ vs.\ $11.1\%$; \PboZeroCallsW{} of matched women place zero outbound calls vs.\ \PboZeroCallsM{} of men), and once real follow-up conversations occur (\PboMatchCallPct{} of matched participants call; \PboMatchDatePct{} date), women's temporary 24-hour surge on looks ($\PboHindLooksJumpW$ points) reverses as intelligence rebounds $\PboHindIntelSnapW$ points ($[+0.8,+6.8]$) and looks recede $\PboHindLooksRecedeW$ points.

\subsection{Generalist Breadth, Staged Capacity, Enforced Delivery, and Empirical Mind Selection}
When a residential band observes every adult across $K$ common domains---foraging, toolmaking, childcare, argument, navigation, and conflict mediation---four theoretical and empirical properties emerge (Fig.~\ref{fig:generalist_verify}):
\begin{enumerate}[leftmargin=1.4em,itemsep=1.5pt,topsep=2pt]
\item \textbf{Breadth aims selection at $g$ and buys time:} Under Corollary~\ref{cor:generalist}, comparing candidates on $K$ shared domains multiplies the selection differential on general capacity relative to any specialty by $K$ (simulated ratios $\PolyRatioOne,\PolyRatioFour,\PolyRatioSixteen$ at $K=1,4,16$), whereas occupational division of labor where each person is judged in their best specialty weakens generalist aim $\PolyDivTen$-fold at $K=10$ (Proposition~\ref{prop:division}). Breadth also makes pre-reproductive assessment feasible: reaching $R^*=0.5$ is impossible for $K\le 5$ domains, takes $\PolyNSix$ events per domain at $K=6$, and requires roughly one event per domain at $K=12$ (Proposition~\ref{prop:time}; Fig.~\ref{fig:generalist_verify}A). Evolutionarily, general-capacity investment dominates total expenditure ($\pi_G>\tfrac12$) once repertoire breadth $K$ exceeds the tissue-to-skill cost ratio $c_G/c_S$, while single-domain exaggeration is taxed as $1/K$ (Proposition~\ref{prop:polymath}, Corollary~\ref{cor:inflation}; Fig.~\ref{fig:generalist_verify}B).
\item \textbf{Staged tests dissociate from unstaged delivery:} Raising the share of announced, staged assessments $d$ from $0.1$ to $0.9$ across $n=\VerN$ occasions lifts capacity validity $R_b$ from $\VerRbLo$ to $\VerRbHi$ while collapsing delivery validity $R_p$ from $\VerRpLo$ to $\VerRpHi$, allowing a candidate who defaults on $25\%$ of everyday obligations to maintain a clean staged record $\VerCleanNine\%$ of the time (Proposition~\ref{prop:staging}; Fig.~\ref{fig:generalist_verify}C).
\item \textbf{Third-party enforcement breaks the take-off trap:} Below the co-evolutionary threshold ($T=\EnfT<1$), neither capacity nor commitment evolves under dyadic choice alone ($e=0$). Enforcing camp-wide sharing of $e=0.2$ or $0.5$ of output~\cite{woodburn1982,boehm1999} immediately lifts equilibrium capacity to $\EnfCapFifth$ and $\EnfCapHalf$ of its maximum while peaking voluntary commitment at partial enforcement ($e^*=\EnfEstar$), whereas full enforcement ($e\to 1$) crowds out selection on internal commitment above threshold ($T=\EnfCrowdT$; Proposition~\ref{prop:enforce}; Fig.~\ref{fig:generalist_verify}D).
\item \textbf{Four-minute SAT blindness versus nine-domain written portfolios:} Pairing speed-dating ratings with participants' objective college SAT scores (\PboSatDatesN{} encounters across \PboSatPartN{} participants~\cite{fisman2006}) and \PboOkcEssaysN{} multi-prompt \texttt{OkCupid-JSE} essays ($N=\PboOkcN$ San Francisco profiles~\cite{kim2015}) tests Propositions~\ref{prop:portfolio},~\ref{prop:time}, and~\ref{prop:twoway} directly (Fig.~\ref{fig:generalist_verify}E,F). One stranger in four minutes explains only \PboSatSingleRTwo{} of a partner's objective SAT ($r=\PboSatSingleR$, 95\% CI $[+0.022,+0.096]$), objective SAT has a null effect on saying ``Yes'' ($r=\PboSatYesR$, $[-0.006,+0.065]$; matched couples $r=\PboSatMatchR$), and even pooling $10+$ four-minute raters stalls at \PboSatPanelRTwo{} ($r=\PboSatPanelR$, $\PboSatPanelGain\times$ single-rater $R^2$ yet trapped at $k=1$). Conversely, across $K=9$ topically distinct \texttt{OkCupid} prompts (within-person cross-prompt engagement $r=+0.467$; $n=\PboOkcDomWritersN$ complete 9-domain writers), single-prompt lexical-sophistication reliability ($r_1=\PboOkcDomROne$) rises $\PboOkcDomGain\times$ to $R_G(9)=\PboOkcDomRNine$. Across writer lexical quintiles ($N=\PboOkcAssortN$), explicit partner requests for curiosity, intelligence, and witty conversation in \texttt{essay9} (``You should message me if\dots'') climb $+40\%$ from \PboOkcAssortLo{} ($Q_1$) to \PboOkcAssortHi{} ($Q_5$; $\PboOkcAssortRatio\times$) and kindness climbs $+38\%$ ($8.3\%\to 11.5\%$), whereas height/money/looks demands stay flat near $2.3\%$ and prompt completion drops sharpest on \texttt{essay8} (``Most private thing I'm willing to admit,'' \PboOkcSelfPct{}$\to$\PboOkcPrivPct{}).
\end{enumerate}

\begin{figure*}[!t]
\centering
\includegraphics[width=\textwidth]{performance-fig-height-attraction.pdf}
\caption{\textbf{Male height percentile versus attraction, the $6'0''$ platform cliff, the straight-versus-gay audience-staging natural experiment, and null completed fertility.} \textbf{(A)}~Male height percentile versus self-rated attraction percentile across $N=\PboWhyteMenN$ Australian men~\cite{whyte2021}: attraction rises from the $\PboWhyteShortAttrPct$th percentile in the shortest bracket to a plateau near the $\PboWhytePeakAttrPct$th at the $\PboWhytePeakHeightPct$th height percentile before dipping to $\PboWhyteTallAttrPct$th at the $\PboWhyteTallHeightPct$th percentile ($\beta_2=\PboWhyteQuadBeta$, \PboWhyteQuadCI). \textbf{(B)}~Reported male stature across $N=\PboHeightMenN$ POV profiles versus Gaussian expectation, exposing a $\PboHeightSixZeroRatio\times$ heap at $6'0''$ ($n=\PboHeightSixZeroN$) and a $-\PboHeightCliffDrop$ cliff at $6'1''$ ($n=\PboHeightSixOneN$). \textbf{(C)}~Audience-targeted staging on the identical San Francisco \texttt{OkCupid-JSE} interface~\cite{kim2015}: straight men ($n=\PboOkcStrMenN$) surge $\PboOkcStrJump$ from $5'11''$ to $6'0''$ before a $-\PboOkcStrCliff$ cliff at $6'1''$, whereas gay men ($n=\PboOkcGayMenN$) show no $6'0''$ surge ($\PboOkcGayJump$), and straight men disclose positive income at \PboOkcIncStrM{}---$\PboOkcIncGayRatio\times$ gay men (\PboOkcIncGayM{}) and $\PboOkcIncFemRatio\times$ straight women (\PboOkcIncStrW{}). \textbf{(D)}~Prime-age POV outcomes (ages 30--49) across stature percentiles: $>3$-year union history rises across the lower-tail floor ($\PboHeightUnionShort\to\PboHeightUnionPeak$) and flattens above median height, while casual-intent share ticks upward in the tall tail ($\PboHeightCasualShort\to\PboHeightCasualTall$). \textbf{(E)}~Cross-population meta-analysis~\cite{lidborg2022}: male height predicts short-term mating ($r=+\PboLidborgHeightMatR$, $k=\PboLidborgHeightMatK$) but shows null association with completed fertility ($r=+\PboLidborgHeightRepR$, \PboLidborgHeightRepCI, $k=\PboLidborgHeightRepK$, $N=\PboLidborgHeightRepN$), whereas functional bodily strength predicts both mating and reproduction ($r=+\PboLidborgBodyRepR$). \textbf{(F)}~Across \PboOkcMsgN{} \texttt{OkCupid-JSE} ``message me if'' essays~\cite{kim2015}, humor (\PboOkcMsgHumor{}), kindness (\PboOkcMsgKind{}), and curiosity/mind (\PboOkcMsgMind{}) outnumber $6'0''$/tall (\PboOkcMsgTall{}) by $\PboOkcHumorTallRatio\times$ and wealth (\PboOkcMsgRich{}) by $\PboOkcHumorRichRatio\times$, while in \PboOkcNoticeN{} ``first noticed'' essays smile (\PboOkcSmilePct{}) and eyes (\PboOkcEyesPct{}) outrank height (\PboOkcHeightNoticePct{}).}\label{fig:height}
\end{figure*}

\begin{table*}[!t]
\caption{Male Stature Percentiles Across Attraction, $6'0''$ Platform Staging, Union History, and Meta-Analytic Fertility.}\label{tab:height}
\scriptsize
\setlength{\tabcolsep}{4.2pt}
\begin{tabularx}{\textwidth}{@{}lrclL@{}}
\toprule
\tabinput{performance-height-table.tex}
\bottomrule
\end{tabularx}
\end{table*}

\subsection{Male Height Percentiles: Threshold Floor, \texorpdfstring{$6'0''$}{6'0"} Staging, and Null Fertility}
Male stature illustrates how a highly legible static cue behaves across perception, strategic online presentation, and completed reproduction (Fig.~\ref{fig:height}; Table~\ref{tab:height}). In Australian microdata ($N=\PboWhyteMenN$ men aged 18--64~\cite{whyte2021}), mapping ordinal stature into cumulative height percentiles reveals a concave floor-and-plateau rather than runaway linearity: mean attraction percentile climbs steeply out of the shortest decile ($\PboWhyteShortAttrPct$th percentile) to $\PboWhytePeakAttrPct$th near the $\PboWhytePeakHeightPct$th height percentile, then flattens and dips slightly to $\PboWhyteTallAttrPct$th at the $\PboWhyteTallHeightPct$th percentile (standardized quadratic term $\beta_2=\PboWhyteQuadBeta$, 95\% CI \PboWhyteQuadCI; Fig.~\ref{fig:height}A)~\cite{stulp2012,stulp2013}. Standardized multivariate regressions confirm that health ($+\PboWhyteHealthSlope$ attraction pct/$\SD$) outweighs stature ($+\PboWhyteHeightSlope$), and income ($+\PboWhyteIncomePartSlope\,\mathrm{pp}/\SD$) outweighs stature ($+\PboWhyteHeightPartSlope\,\mathrm{pp}/\SD$) for current partnering.

On digital platforms ($N=\PboHeightMenN$ men with explicit stature on Plenty of Fish), lexical filtering around round-number cutoffs induces Campbell--Goodhart gaming~\cite{campbell1979} (Fig.~\ref{fig:height}B): profiles reporting $6'0''$ ($183\,\mathrm{cm}$, $n=\PboHeightSixZeroN$) reach $\PboHeightSixZeroRatio\times$ Gaussian expectation before plunging $-\PboHeightCliffDrop$ at $6'1''$ ($n=\PboHeightSixOneN$). Comparing straight ($n=\PboOkcStrMenN$) and gay ($n=\PboOkcGayMenN$) men on the identical San Francisco \texttt{OkCupid-JSE} interface~\cite{kim2015} proves that this cliff reflects \emph{audience-targeted staging} rather than generic measurement rounding (Fig.~\ref{fig:height}C): straight men surge $\PboOkcStrJump$ from $5'11''$ ($3{,}893$) to $6'0''$ ($4{,}559$) before a $-\PboOkcStrCliff$ cliff at $6'1''$ ($2{,}427$), whereas gay men exhibit \emph{no} $5'11''\to 6'0''$ surge ($471\to 465$, $\PboOkcGayJump$), and straight men simultaneously disclose positive income at \PboOkcIncStrM{}---$\PboOkcIncGayRatio\times$ gay men (\PboOkcIncGayM{}) and $\PboOkcIncFemRatio\times$ straight women (\PboOkcIncStrW{}, $n=\PboOkcStrWomenN$). Across ages 30--49 on POV, prior $>3$-year union attainment exhibits a lower-tail floor ($\PboHeightUnionShort$ at $5'1''$ rising to $\PboHeightUnionPeak$ at $6'1''$ and $\PboHeightUnionTall$ at $6'6''$), while casual-intent share ticks upward in the tall tail ($\PboHeightCasualShort\to\PboHeightCasualTall$; Fig.~\ref{fig:height}D). Synthesizing $k=\PboLidborgHeightRepK$ populations ($N=\PboLidborgHeightRepN$) in Lidborg et al.~\cite{lidborg2022} shows that while male height positively predicts mating access ($r=+\PboLidborgHeightMatR$, \PboLidborgHeightMatCI, $k=\PboLidborgHeightMatK$), its net correlation with completed fertility is statistically indistinguishable from zero ($r=+\PboLidborgHeightRepR$, \PboLidborgHeightRepCI), whereas functional muscular strength predicts both mating and reproduction ($r=+\PboLidborgBodyRepR$, \PboLidborgBodyRepCI; Fig.~\ref{fig:height}E; Proposition~\ref{prop:aim}). Meanwhile, across \PboOkcMsgN{} \texttt{OkCupid} ``You should message me if\dots'' essays~\cite{kim2015}, explicit invitations for humor (\PboOkcMsgHumor{}), kindness (\PboOkcMsgKind{}), and curiosity/mind (\PboOkcMsgMind{}) outnumber $6'0''$/height cutoffs (\PboOkcMsgTall{}) by $10.0\times$--$\PboOkcHumorTallRatio\times$ and wealth (\PboOkcMsgRich{}) by $\PboOkcHumorRichRatio\times$ (Fig.~\ref{fig:height}F).

\begin{figure*}[!t]
\centering
\includegraphics[width=\textwidth]{performance-fig-ethnography.pdf}
\caption{\textbf{Ethnographic marriage juries, constraint versus exit, male waiting times, and effective ancestor ratios.} \textbf{(A)}~Who selects a spouse across the Standard Cross-Cultural Sample (SCCS, 186 societies) and hunter-gatherer subset~\cite{murdock1969,kirby2016}: kin and elders participate in the majority of unions, and autonomous choice is rarer for daughters (\CcAloneW\% overall, \CcAloneWFor\% in foragers) than sons (\CcAloneM\%, \CcAloneMFor\%). \textbf{(B)}~Prevalence of marriage constraint and coercion alongside symmetric divorce access (\CcDivEqPct\%) and frequent exit (\CcDivCommonPct\%) in coded societies. \textbf{(C)}~Mean age at first marriage across \CcAgeN{} hunter-gatherer groups~\cite{binford2001}: men marry $\CcGap\,\mathrm{yr}$ later than women ($20.8$ vs.\ $14.1\,\mathrm{yr}$), creating a multi-year brideservice and foraging audit window. \textbf{(D)}~Female-to-male variance effective population ratio $N_{e,f}/N_{e,m}=[(1-q_f)/(1-q_m)][(1+\alpha_m)/(1+\alpha_f)]$ (Proposition~\ref{prop:ancestors}): a uniparental genomic ratio of $2.0$ does not imply half of men were excluded, and completed U.S.\ cohorts (1900--1979, $n=\VgAncN$) sit near $\VgNeRatio$.}\label{fig:ethnography}
\end{figure*}

\begin{figure*}[!t]
\centering
\includegraphics[width=\textwidth]{performance-fig-world-synthesis.pdf}
\caption{\textbf{Cross-national partner priorities, shared national culture, and global digital platform gradients.} \textbf{(A)}~Nine-trait stated partner importance dumbbell across $N=\PboWhyteN$ Australians~\cite{whyte2021}: trust ($\PboWhyteTrustW,\PboWhyteTrustM/100$), emotional connection ($\PboWhyteEmotConnW,\PboWhyteEmotConnM$), openness, and intelligence ($\PboWhyteIntelW,\PboWhyteIntelM$) outrank physique and income in both sexes. \textbf{(B)}~Within-country sex differences (women minus men, points of 30) across \PboKowalCountryN{} countries ($n=\PboKowalBudgetN$~\cite{kowal2025}), showing shared directional tilts alongside wide national spread. \textbf{(C)}~Cross-national co-variation of men's and women's 30-point mate budgets: national culture shifts both sexes in lockstep ($r=\PboKowalRelCorr$ religion, $r=\PboKowalKindCorr$ kindness). \textbf{(D)}~Union-entry partnering slopes per self-rated $\SD$ across $\PboDoorN$ adults~\cite{kowal2025}: financial resources ($\PboDoorMoneyM\,\mathrm{pp}$) and looks ($\PboDoorLooksM\,\mathrm{pp}$) predict men's union entry, whereas self-rated kindness is near zero ($\PboDoorKindM,\PboDoorKindW\,\mathrm{pp}$). \textbf{(E)}~Within-gender HON beauty-percentile drop from ages 18--29 to 40--59 across \PboHonCountriesStrictAge{} strict countries ($n\ge 500$ per sex): steeper for women in \PboHonFemaleSteeperCountries/\PboHonCountriesStrictAge{} nations (median $-\PboHonAgeDropMedW$ vs.\ $-\PboHonAgeDropMedM\,\mathrm{pp}$). \textbf{(F)}~National platform participation scale versus female-minus-male like-rate advantage across \PboHonCountriesFiveHundred{} countries ($r=\PboHonPartLikeCorr$), reflecting upload selectivity in smaller markets.}\label{fig:world}
\end{figure*}

\begin{figure*}[!t]
\centering
\includegraphics[width=\textwidth]{performance-fig-love-bond.pdf}
\caption{\textbf{Primary pair-bond devotion versus secondary entry filters across 175 countries, 43 longitudinal couples studies, Stanford \texttt{HCMST}, and 1.05M essays (Proposition~\ref{prop:love}).} \textbf{(A)}~Country-FE relationship-quality horse race ($N=\PboChkCouplesN$ couples, \PboChkCountriesN{} countries~\cite{kowal2025,conroybeam2016}): partner kindness ($\PboChkKindSD\,\SD$) outruns 5-trait checklist fit ($\PboChkFitSD\,\SD$) by $\PboChkKindFitRatio\times$, partner money ($\PboChkResSD\,\SD$), and ``marrying up'' ($\PboChkMarryUpSD\,\SD$) across all four cultural regimes. \textbf{(B)}~Random-forest synthesis across \PboJoelDatasetsN{} longitudinal couples datasets ($11{,}196$ couples~\cite{joel2020}): perceived commitment and dyadic conduct explain \PboJoelRelBaseRTwo{} of baseline satisfaction (\PboJoelRelFollowRTwo{} follow-up), $\PboJoelBaseRatio\times$ static partner traits (\PboJoelTraitBaseRTwo{}). \textbf{(C)}~Stanford \texttt{HCMST} ($N=\PboHcmstN$~\cite{rosenfeld2019}): apps compress pre-romance acquaintance $\PboHcmstAuditRatio\times$ ($\PboHcmstAppMeanYr$ vs.\ $\PboHcmstInstMeanYr\,\mathrm{yr}$ institutional/kin), while surviving intact unions show equal quality ($\PboHcmstQualApp$ vs.\ $\PboHcmstQualInst/5$). \textbf{(D)}~25-year commitment arc ($N=\PboLoveArcN$~\cite{sternberg1986,kowal2025}; $\PboLoveCommStart\to\PboLoveCommEnd/5$), female emotional-over-sexual jealousy ($\PboJealEmoW$ vs.\ $\PboJealSexW/7$~\cite{buss1989}), and observational social-media infidelity slope ($\PboSocMedInfSlope\,\mathrm{pp/h}$). \textbf{(E)}~Self-written themes across \PboVoiceEssaysN{} POV essays (\PboVoiceFamilyW{} W write family/home, $\PboVoiceFamilyTallRatio\times$ $6'0''$ mentions). \textbf{(F)}~Coded divorce causes across \PboBetzigSocN{} SCCS societies~\cite{betzig1989}: betrayal (\PboBetzigAdulteryPct) and cruelty (\PboBetzigCrueltyPct) occur at $\PboBetzigRatio\times$ aging looks (\PboBetzigAgingPct).}\label{fig:love_bond}
\end{figure*}

\begin{figure*}[!t]
\centering
\includegraphics[width=\textwidth]{performance-fig-primary-secondary-dynamics.pdf}
\caption{\textbf{Cross-source male--female mating dynamics separating secondary sorting from primary devotion (Proposition~\ref{prop:lifecycle}).} \textbf{(A)}~Columbia speed-dating cross-sex mind-reading error ($N=\SdN$~\cite{fisman2006}): both sexes overestimate opposite-sex looks weight ($\PboMindReadLooksW$ W, $\PboMindReadLooksM$ M pts) and underestimate intelligence ($\PboMindReadIntelW,\PboMindReadIntelM$ pts), while men place $\PboCallRatioMF\times$ more post-match calls (\PboZeroCallsW{} W zero calls). \textbf{(B)}~The homophily paradox across $149$--$150$ countries ($N=47.4\mathrm{k}$--$50.0\mathrm{k}$~\cite{kowal2025}): couples sort strongly on politics ($r=\PboHomoPolR$), social class ($r=\PboHomoClsR$), and education ($r=\PboHomoEduR$), yet primary kindness ($+0.39\,\SD$) beats secondary homophily slopes ($\PboHomoPolSlope,\PboHomoClsSlope,\PboHomoEduSlope\,\SD$) by $\PboHomoPolRatio\text{--}\PboHomoEduRatio\times$. \textbf{(C)}~Country-FE partner reverence versus ``settling'' spline ($N=50{,}180$~\cite{kowal2025}): settling below self on kindness penalizes love $\PboSettleRevRatioW\times$ (W) and $\PboSettleRevRatioM\times$ (M) more steeply than surplus adds, while money reverence is null and mutual looks reverence cuts infidelity $-\PboRevLooksInfDropW/-\PboRevLooksInfDropM\,\mathrm{pp}$. \textbf{(D)}~Daily appearance/physical courtship staging peaks in dating (\PboEffortDatingW/\PboEffortDatingM\,min/d) and drops $-\PboEffortDropPctW/-\PboEffortDropPctM$ in marriage ($N=65{,}916$ observational), alongside widening male-over-female casual \texttt{SOI} desire ($\PboSoiRatioYoung\to\PboSoiRatioOld\times$). \textbf{(E)}~Stanford \texttt{HCMST} ($N=2{,}578$ ever-cohabited~\cite{rosenfeld2019}): rushing from romance onset to cohabitation in $\le 3$ months triples breakup (\PboCohabRushPct{} vs.\ \PboCohabSlowPct{} at $2.5+$\,yr, $\PboCohabRushRatio\times$), while marital exits flip to \PboDivUnilatW{} female unilateral initiation (\PboDivHusRepW{} in husbands' reports vs.\ \PboNonMarUnilatW{} non-marital). \textbf{(F)}~Across \PboVoiceEssaysN{} POV essays (ages 18--79), hearth cooking verbs surge $\PboVerbCookRiseW\times$ (W) and $\PboVerbCookRiseM\times$ (M) as party verbs fall $-\PboVerbPartyDrop$, while 40--49 divorced women cut marriage intent $-\PboDivMarryDropW$ (keeping \PboDivSeriousMidW{} serious bond) and widowed men exceed widowed women.}\label{fig:dynamics}
\end{figure*}

\section{Primary Versus Secondary Mating Criteria Across Cultures and Long-Term Unions}\label{sec:cultures}

\subsection{Kin Juries, Forager Waiting Windows, and Effective Ancestors}
In small-scale societies, partner selection is rarely an unmediated dyadic transaction between strangers (Fig.~\ref{fig:ethnography}). Across the Standard Cross-Cultural Sample (SCCS, 186 societies~\cite{murdock1969,kirby2016}), daughters choose spouses autonomously without parental approval in only \CcAloneW\% of societies (\CcAloneWFor\% among hunter-gatherers), compared with \CcAloneM\% for sons (\CcAloneMFor\% among hunter-gatherers), while parents and kin arrange or co-approve the overwhelming majority (Fig.~\ref{fig:ethnography}A)~\cite{apostolou2007,walker2011}. Ethnographic records also document substantial coercion and constraint---bride consent is non-essential in 67\% of 57 coded societies, war captivity of women occurs in \CcCapPct\% of \CcCapN{}, and wife-beating in \CcBeatPct\% of \CcBeatN{}---which dilutes assessment-based selection linearly by $(1-\kappa)$ while redirecting $\kappa$ toward force (Eq.~\eqref{eq:coercion}). Yet exit remains widespread: divorce is equally accessible to both spouses in \CcDivEqPct\% of coded societies (\PboCcDivEqFor\% in foragers; Fig.~\ref{fig:ethnography}B)~\cite{blurtonjones2000,scelza2013}.

Across \CcAgeN{} hunter-gatherer groups in Binford's compilation~\cite{binford2001}, first marriage occurs at a mean age of $\CcAgeM\,\mathrm{yr}$ for men versus $\CcAgeW\,\mathrm{yr}$ for women---a $+\CcGap\,\mathrm{yr}$ male waiting window (95\% CI \CcGapCI\,yr; men older in \CcMenOlderPct\% of groups; Fig.~\ref{fig:ethnography}C). That multi-year interval, reinforced by brideservice in \CcBrideService\% of societies and bridewealth in \CcBridewealth\%, places prospective husbands under continuous unstaged observation ($(1-d)n\gg 1$) by parents-in-law and campmates before reproductive commitment stabilizes~\cite{collier1981,marlowe2004}. Proposition~\ref{prop:ancestors} also clarifies a widely misinterpreted genomic statistic: Y-chromosome versus mitochondrial effective-size ratios near $N_{e,f}/N_{e,m}\approx 2$~\cite{wilder2004,karmin2015} do not imply that half of ancestral men left zero descendants; identical ratios arise from modest father-to-father family-size dispersion $\alpha_m$ or patrilineal kin-group drift~\cite{zeng2018}, while completed U.S.\ cohorts (1900--1979, $n=\VgAncN$) sit near $\VgNeRatio$ (\VgNeRatioCI; Fig.~\ref{fig:ethnography}D).

\subsection{Stated Trust, Shared National Culture, and Cross-National Platform Gradients}
Modern cross-national microdata confirm that when people evaluate long-run partners rather than four-minute strangers, delivery, emotional connection, and mind sit at the top of the hierarchy (Fig.~\ref{fig:world}). Across $N=\PboWhyteN$ Australians~\cite{whyte2021}, interpersonal trust ($\PboWhyteTrustW/100$ women, $\PboWhyteTrustM/100$ men), emotional connection ($\PboWhyteEmotConnW,\PboWhyteEmotConnM$), openness, and intelligence ($\PboWhyteIntelW,\PboWhyteIntelM$) outrank physical build and income in both sexes (Fig.~\ref{fig:world}A). Across \PboKowalBudgetN{} constrained 30-point mate budgets from \PboKowalCountryN{} countries in Kowal et al.~\cite{kowal2025} (verified across \PboKowalStdCountryN{} age-standardized nations), women allocate slightly more to financial prospects and kindness while men allocate more to looks (Fig.~\ref{fig:world}B), yet national institutional culture shifts men's and women's priorities almost in unison: country-mean allocations correlate across sexes at $r=\PboKowalRelCorr$ for religion and $r=\PboKowalKindCorr$ for kindness, dwarfing within-country sex gaps (Fig.~\ref{fig:world}C; assortative partner matching likewise peaks on shared religion $r=\PboAssortReligR$ and kindness $r=\PboAssortKindR$). At the union-entry margin across $\PboDoorN$ adults~\cite{kowal2025} and 45--70-country surveys ($n=\WdN$~\cite{walter2020}), men's self-rated financial standing ($\PboDoorMoneyM\,\mathrm{pp}/\SD$) and looks ($\PboDoorLooksM\,\mathrm{pp}$) strongly predict having a partner, whereas self-rated kindness barely moves entry ($\PboDoorKindM\,\mathrm{pp}$ in men, $\PboDoorKindW\,\mathrm{pp}$ in women; Fig.~\ref{fig:world}D).

Comparing \PboHonCountriesFiveHundred{} countries with $\ge 500$ Hot-or-Not profiles per sex (\PboHonCountriesHundred{} countries at $\ge 100$ per sex) reveals two global regularities of single-photograph evaluation. First, among \PboHonCountriesStrictAge{} countries meeting strict age-cell thresholds, within-gender beauty percentiles fall more steeply across adulthood for women than men in \PboHonFemaleSteeperCountries/\PboHonCountriesStrictAge{} nations (median $-\PboHonAgeDropMedW$ vs.\ $-\PboHonAgeDropMedM\,\mathrm{pp}$; Fig.~\ref{fig:world}E; Corollary~\ref{cor:cue}). Second, log national platform participation correlates negatively with the female-minus-male like-rate gap ($r=\PboHonPartLikeCorr$; Fig.~\ref{fig:world}F), showing that thin-market adopter self-selection inflates observed ratings unless upload selectivity is audited.

\subsection{Primary Devotion Over Secondary Homophily, the Settling Kink, Rushed Lock-In, and Asymmetric Marital Exit}

\begin{table*}[!t]
\caption{Taxonomy and Empirical Synthesis of Secondary Entry Criteria Versus Primary Pair-Bond Criteria Across Fifteen Datasets.}\label{tab:primary_secondary}
\scriptsize
\renewcommand{\arraystretch}{0.83}
\setlength{\tabcolsep}{3.8pt}
\begin{tabularx}{\textwidth}{@{}>{\raggedright\arraybackslash}p{0.18\textwidth}>{\raggedright\arraybackslash}p{0.18\textwidth}>{\raggedright\arraybackslash}p{0.185\textwidth}LL@{}}
\toprule
Mating Criterion & Empirical Scale & Stranger Legibility ($\rho_0$) & Entry / Assortative Sorting Effect & Realized Pair-Bond, Divorce \& Fertility Return \\
\midrule
\tabinput{performance-primary-secondary-table.tex}
\bottomrule
\end{tabularx}
\vspace{-3.5pt}
\tablenote{Secondary criteria govern low-acquaintance market filtering and assortative sorting ($\rho_X$ high) while contributing near zero inside established unions ($\alpha_X\approx 0$); primary criteria are noisy in single snapshots ($\rho_{K,0}\ll\rho_X$) yet dominate long-term satisfaction, fidelity, and juvenile provisioning ($\alpha_K\gg|\alpha_X|$; Propositions~\ref{prop:love}--\ref{prop:lifecycle}).}
\end{table*}

Why do searchable dating markets look obsessed with income, youth, $6'0''$ cutoffs, and political or educational credentials when lived marriages revolve around kindness, mutual reverence, and commitment? Propositions~\ref{prop:love} and~\ref{prop:lifecycle} unify these facts by separating \emph{secondary entry criteria} ($\rho_X$ high, $\alpha_X\approx 0$) from \emph{primary pair-bond criteria} ($\rho_{K,0}\ll\rho_X$, $\alpha_K,\alpha_D\gg|\alpha_X|$; Table~\ref{tab:primary_secondary}; Figs.~\ref{fig:love_bond}--\ref{fig:dynamics}):
\begin{enumerate}[leftmargin=1.4em,itemsep=1.5pt,topsep=2pt]
\item \textbf{Kindness outruns checklist fit, ``marrying up,'' and demographic homophily:} Across $N=\PboBondPrqcN$ partnered adults in 175 countries~\cite{kowal2025}, financial resources predict men's union entry ($\PboDoorMoneyM\,\mathrm{pp}/\SD$) yet collapse near zero inside unions ($\PboBondMoneyW,\PboBondMoneyM\,\SD$; infidelity $\PboBondInfMoneyW\,\mathrm{pp}$, $N=\PboBondInfN$), whereas partner kindness dominates love ($\PboBondKindW,\PboBondKindM\,\SD$) and cuts infidelity ($\PboBondInfKindW,\PboBondInfKindM\,\mathrm{pp}/\SD$; women even rank kindness $\#1$ for short-term mates, $\PboIdealStKindW/11$ vs.\ $\PboIdealStLooksW$ looks). In a country-FE horse race ($N=\PboChkCouplesN$ couples across \PboChkCountriesN{} countries; Fig.~\ref{fig:love_bond}A), kindness ($\PboChkKindSD\,\SD$) beats Conroy-Beam 5-trait idiographic checklist fit ($\PboChkFitSD\,\SD$~\cite{conroybeam2016}) by $\PboChkKindFitRatio\times$, partner money ($\PboChkResSD\,\SD$) by $8.5\times$, and ``marrying out of your league'' ($\PboChkMarryUpSD\,\SD$, $[-0.03,-0.01]$; $\PboChkMarryUpFiveSD\,\SD$ 5-trait controlled) across egalitarian, low-equality, individualist, and collectivist regimes. Crucially, across $149$--$150$ countries ($N=47{,}410$--$49{,}984$), couples sort heavily on secondary social labels---political ideology ($r=\PboHomoPolR$, $66.7\%$ exact match), social class ($r=\PboHomoClsR$, $65.5\%$), and education ($r=\PboHomoEduR$, $52.0\%$)---yet inside-union country-FE quality slopes for similarity are tiny ($\PboHomoPolSlope$, $\PboHomoClsSlope$, and $\PboHomoEduSlope\,\SD$), so primary partner kindness beats secondary demographic homophily by $\PboHomoPolRatio\times$, $\PboHomoClsRatio\times$, and $\PboHomoEduRatio\times$ (Fig.~\ref{fig:dynamics}B).
\item \textbf{Partner reverence versus the ``settling'' kink, and 43 longitudinal couples studies:} Estimating piecewise country-FE splines in $(\text{mate}-\text{self})$ ratings controlling for own trait ($N=50{,}180$; Fig.~\ref{fig:dynamics}C) explains why mean ``marrying up'' looks flat on money while partner admiration anchors love: settling below oneself on primary kindness penalizes love by $\PboSettleKindSlopeW\,\SD/\mathrm{pt}$ in women and $\PboSettleKindSlopeM\,\SD/\mathrm{pt}$ in men---$\PboSettleRevRatioW\times$ and $\PboSettleRevRatioM\times$ steeper than positive reverence surplus ($\PboRevKindSlopeW$ and $\PboRevKindSlopeM\,\SD/\mathrm{pt}$)---whereas reverence surplus on financial resources is null ($\PboRevResSlopeW$ W, $\PboRevResSlopeM$ M), and perceiving one's partner as more physically attractive than oneself cuts reported infidelity by $-\PboRevLooksInfDropW\,\mathrm{pp}$ in women and $-\PboRevLooksInfDropM\,\mathrm{pp}$ in men. Independently, across \PboJoelDatasetsN{} longitudinal couples datasets ($11{,}196$ couples~\cite{joel2020}), perceived partner commitment and dyadic responsiveness explain \PboJoelRelBaseRTwo{} of baseline satisfaction (\PboJoelRelFollowRTwo{} at 1--4\,yr follow-up), $\PboJoelBaseRatio\times$ all static partner traits combined (\PboJoelTraitBaseRTwo{}, \PboJoelTraitFollowRTwo{}; Fig.~\ref{fig:love_bond}B).
\item \textbf{Courtship staging drop, app witness compression, and the $\le 3$-month cohabitation rush hazard:} Consistent with Proposition~\ref{prop:lifecycle}(i) (\texttt{E23} closed form $-\PboETwentyThreeDropPct$), daily appearance and physical effort across $N=65{,}916$ adults~\cite{kowal2025} peaks during active dating (\PboEffortDatingW\,min/d women, \PboEffortDatingM\,min/d men) and drops $-\PboEffortDropPctW$ (to \PboEffortMarriedW\,min/d) in married women and $-\PboEffortDropPctM$ (to \PboEffortMarriedM\,min/d) in married men (observational cross-section; Fig.~\ref{fig:dynamics}D). In Stanford \texttt{HCMST} ($N=\PboHcmstN$, \PboHcmstCouplesN{} meeting-coded couples~\cite{rosenfeld2019}), dating apps compress pre-romance acquaintance $\PboHcmstAuditRatio\times$ ($\PboHcmstAppMeanYr\,\mathrm{yr}$, median $\PboHcmstAppMedYr$, \PboHcmstAppZeroPct{} zero years, \PboHcmstAppZeroTiePct{} zero prior ties vs.\ $\PboHcmstInstMeanYr\,\mathrm{yr}$ institutional/kin and $\PboHcmstSchoolMeanYr\,\mathrm{yr}$ school, $\PboHcmstSchoolRatio\times$; surviving intact-union quality null $\PboHcmstQualApp$ vs.\ $\PboHcmstQualInst/5$, adjusted $\PboHcmstQualAdjDiff$; Fig.~\ref{fig:love_bond}C). Directly confirming Proposition~\ref{prop:lifecycle}(ii), among $2{,}578$ ever-cohabited \texttt{HCMST} couples, rushing from romance onset into cohabitation within $\le 3$ months triples subsequent union dissolution (\PboCohabRushPct{} vs.\ \PboCohabSlowPct{} when waiting $2.5+$ years, $\PboCohabRushRatio\times$; cohort- and age-adjusted penalty $\PboCohabRushAdjPP\,\mathrm{pp}$, $[+1.4,+6.6]$; Fig.~\ref{fig:dynamics}E).
\item \textbf{The 25-year commitment arc, voice, and marriage-specific female unilateral exit:} Across \PboLoveArcN{} couples~\cite{sternberg1986,kowal2025}, passion moderates ($\PboLovePassStart\to\PboLovePassEnd/5$) while intimacy stays high (\PboLoveIntimEnd) and commitment rises ($\PboLoveCommStart\to\PboLoveCommEnd/5$; love documented in \PboJankowiakPct{} of 166 societies~\cite{jankowiak1992}; marry-without-love median $\PboMarryNoLoveMedW,\PboMarryNoLoveMedM/100$, $N=\PboMarryNoLoveN$~\cite{frank1988}), women report stronger emotional than sexual jealousy ($\PboJealEmoW$ vs.\ $\PboJealSexW/7$, $\PboJealDiffW$~\cite{buss1989}), and daily social-media browsing tracks higher perceived outside options ($\PboSocMedQoaSlope\,\SD/\mathrm{h}$), lower quality ($\PboSocMedPrqcSlope\,\SD/\mathrm{h}$), and higher infidelity ($\PboSocMedInfSlope\,\mathrm{pp/h}$, \PboSocMedLowInf{}$\to$\PboSocMedHighInf{}; Fig.~\ref{fig:love_bond}D~\cite{rusbult1998}). In \PboVoiceEssaysN{} self-written POV essays (\PboVoiceEssaysW{} W, \PboVoiceEssaysM{} M), women invoke family/home (\PboVoiceFamilyW; men \PboVoiceFamilyM), laughter (\PboVoiceLaughW,\PboVoiceLaughM), kindness (\PboVoiceKindW,\PboVoiceKindM), and loyalty (\PboVoiceHonestW,\PboVoiceHonestM) at $\PboVoiceHonestTallRatio\times$--$\PboVoiceFamilyTallRatio\times$ $6'0''$/tall mentions (\PboVoiceTallW) and $\PboVoiceKindRichRatio\times$ wealth (\PboVoiceRichW; Fig.~\ref{fig:love_bond}E). When primary devotion fails under legal and domestic lock-in, exit turns sharply asymmetric (Proposition~\ref{prop:lifecycle}(iii)): in \texttt{HCMST}, non-marital unilateral breakups ($n=286$) are near-symmetric (\PboNonMarUnilatW{} female; $50.0\%$ in male reports), whereas marital divorces ($n=67$) flip to \PboDivUnilatW{} female unilateral initiation (\PboDivHusRepW{} in husbands' own reports; \texttt{E23} \PboETwentyThreeDivW; Fig.~\ref{fig:dynamics}E), paralleling \PboBetzigSocN{} SCCS societies~\cite{betzig1989} where betrayal (\PboBetzigAdulteryPct) and cruelty (\PboBetzigCrueltyPct) trigger divorce at $\PboBetzigRatio\times$ aging looks (\PboBetzigAgingPct; Fig.~\ref{fig:love_bond}F).
\end{enumerate}

\begin{figure*}[!t]
\centering
\includegraphics[width=\textwidth]{performance-fig-hon-architecture.pdf}
\caption{\textbf{Attention concentration and portfolio decoupling across \PboHonRows{} Hot-or-Not profiles (\PboHonEligible{} with $\ge 20$ votes).} \textbf{(A)}~Lorenz concentration of votes and likes: women receive $\PboHonAttentionRatio\times$ median votes ($\PboHonVotesMedW$ vs.\ $\PboHonVotesMedM$), while top-decile profiles capture \PboHonTopDecileLikesW{} (women, Gini \PboHonLikesGiniW) and \PboHonTopDecileLikesM{} (men, Gini \PboHonLikesGiniM) of all positive likes. \textbf{(B)}~Multi-photo upload portfolio (0 to 5+ inventoried photographs): extra photos raise median vote volume $+50\%$ without lifting mean beauty percentile above $\sim\!50$th (within-gender age trajectories appear in Fig.~\ref{fig:pov}B).}\label{fig:hon_arch}
\end{figure*}

\begin{figure*}[!t]
\centering
\includegraphics[width=\textwidth]{performance-fig-holoface-ceiling.pdf}
\caption{\textbf{The \texttt{HoloFace-L1} 43-wave deep-learning benchmark, exact 4-family Shapley attribution, and finite-voter psychometric ceiling ($10{,}000$ profiles, $20{,}000$ portraits).} \textbf{(A)}~Held-out test progression from raw UV baseline H1 ($\rho=0.234$, MAE $0.4831$) through strict zero-metadata \texttt{HoloFace-Pure} ($\rho=\PboHoloPureSpearman$, $\rho^2=\PboHoloPureRankVar$, MAE $\PboHoloPureTestMAE$) to multimodal \texttt{HoloFace-Full} (\texttt{HoloFace-L1}: Test MAE $\PboHoloFullTestMAE$, single-image MAE $\PboHoloFullSingleMAE$, $\rho=\PboHoloFullSpearman$, $\rho^2=\PboHoloFullRankVar$). \textbf{(B)}~Exact $2^4=16$ Shapley decomposition of zero-metadata \texttt{HoloFace-Pure} rank variance: $\PboShapAnatomyTotal$ stems from 12-space foundation appearance ($\PboShapFoundShare$), 3D Kendall Riemannian morphology ($\PboShapGeomShare$), and skin chromophore/spectral albedo ($\PboShapSkinShare$), versus only $\PboShapPoseShare$ from changeable pose, smile, and framing. \textbf{(C)}~Finite-voter psychometric Oracle ceiling ($\hat\sigma_{\text{voter}}^2=\PboVoterSigmaSq$ on $1$--$10$ rating scale, fitted strictly on $8{,}000$ training OOF residuals) versus \texttt{HoloFace-Full} and \texttt{HoloFace-Pure} across vote-depth tiers (20--49, 50--199, 200+, 500+), isolating the $\PboIrreducibleGapFiveHundred$ (500+ votes, Oracle $\rho=\PboOracleFiveHundredR$) to $\PboIrreducibleGapAll$ (all test, Oracle $\rho=\PboOracleAllR$) irreducible single-context consensus gap. \textbf{(D)}~Skew-unbiased $L_2$ calibration eliminating male rating bias ($+0.088\to\PboHoloMaleLTwoBias$, RMSE $\PboHoloLTwoRMSE$) alongside the $k=63$ human MHC chemosensory meta-analytic null~\cite{havlicek2020}.}\label{fig:holoface}
\end{figure*}

\section{Digital Platforms at Scale, the \texorpdfstring{\texttt{HoloFace-L1}}{HoloFace-L1} Ceiling, Facial Attribution, and Completed Fertility}\label{sec:platforms}

\subsection{Platform Attention Concentration and Portfolio Decoupling}
Auditing all \PboHonRows{} canonical Hot-or-Not records (\PboHonEligible{} eligible profiles aged 18--79 with $\ge 20$ valid votes and \PboHonFiles{} inventoried images; Appendix Table~\ref{tab:hon_sens}) exposes three structural properties of single-context crowd evaluation (Fig.~\ref{fig:hon_arch}; life-course age curve in Fig.~\ref{fig:pov}B). First, within-gender beauty rank declines much more steeply across adulthood for women than men: women's mean percentile drops $-27.8$ points from $\PboBeautyYoungW$ at ages 25--29 (top-decile share \PboBeautyTopYoungW) to $\PboBeautyOlderW$ at ages 45--49 (\PboBeautyTopOlderW), compared with a gentler $-20.1$-point decline in men from $\PboBeautyYoungM$ to $\PboBeautyOlderM$ after an early-twenties peak (Fig.~\ref{fig:pov}B). Second, voter attention and positive evaluations are heavily concentrated: female profiles receive $\PboHonAttentionRatio\times$ as many median votes as male profiles ($\PboHonVotesMedW$ vs.\ $\PboHonVotesMedM$), and the top decile captures \PboHonTopDecileLikesW{} of all female likes (Gini \PboHonLikesGiniW) and \PboHonTopDecileLikesM{} of all male likes (Gini \PboHonLikesGiniM; Fig.~\ref{fig:hon_arch}A). Third, uploading additional photographs (1 to 5+ images) raises median vote volume $+50\%$ without shifting mean beauty percentile above $\sim\!50$th (Fig.~\ref{fig:hon_arch}B): repeating static snapshots within the same staged presentation genre raises exposure ($m$) without adding multi-context trait validity ($k$).

\subsection{The \texorpdfstring{\texttt{HoloFace-L1}}{HoloFace-L1} 43-Wave Benchmark and Psychometric Voter Ceiling}
How much of human consensus rating can be predicted from static portraits and profile metadata, and how much of the residual reflects finite voter noise versus an irreducible single-context ceiling? Across our 43-wave benchmark on $10{,}000$ profiles and $20{,}000$ portraits (8,000 train / 1,000 validation / 1,000 held-out test, zero identity overlap; Fig.~\ref{fig:holoface}; Appendix Table~\ref{tab:holoface}), progressive nuisance removal and multimodal integration establish four quantitative bounds:
\begin{enumerate}[leftmargin=1.4em,itemsep=1.5pt,topsep=2pt]
\item \textbf{Pure-face predictability and 4-family Shapley anatomy:} Stripping all demographic and tabular metadata, \texttt{HoloFace-Pure} lifts held-out pure-face Spearman correlation to $\rho=\PboHoloPureSpearman$ ($\rho^2=\PboHoloPureRankVar$, Test MAE $\PboHoloPureTestMAE$; rising to $\rho=\PboHoloPureFiveHundredR$, $\rho^2=49.5\%$ among well-voted 500+-vote profiles; Fig.~\ref{fig:holoface}A). Exact Shapley attribution across all $2^4=16$ visual family subsets proves that $\PboShapAnatomyTotal$ of pure-face rank variance derives from 12-space foundation and cortical appearance ($\PboShapFoundShare$), 3D craniofacial and Kendall Riemannian landmark geometry without pixel intensities ($\PboShapGeomShare$), and skin chromophore/spectral micro-relief ($\PboShapSkinShare$), whereas changeable camera pose, smile curvature, and framing contribute only $\PboShapPoseShare$ (Fig.~\ref{fig:holoface}B).
\item \textbf{Multimodal \texttt{HoloFace-Full} optimum and skew calibration:} Combining 29 out-of-fold visual, tabular, and cross-modal predictors via validation cross-validated exact-$L_1$ linear programming achieves a new benchmark best across all 43 waves: Validation MAE $\PboHoloFullValMAE$, Candidate Test MAE $\PboHoloFullTestMAE$, Single-Image Test MAE $\PboHoloFullSingleMAE$, within-gender percentile MAE $\PboHoloFullPctMAE$ points, and Test Spearman $\rho=\PboHoloFullSpearman$ ($\rho^2=\PboHoloFullRankVar$; $\rho=\PboHoloFullFiveHundredR$ at 500+ votes). Because male ratings are right-skewed while median-optimal $L_1$ loss targets conditional medians, pairing the $L_1$ head with a ridge $L_2$ calibration head achieves Test RMSE $\PboHoloLTwoRMSE$ while reducing male mean bias from $+0.088$ to $\PboHoloMaleLTwoBias$ (Fig.~\ref{fig:holoface}D).
\item \textbf{Deattenuated single-context consensus gap:} Fitting heteroscedastic squared OOF residuals $(y_i-\hat f_{\text{oof}}(x_i))^2=\sigma_{\text{gap}}^2+\sigma_{\text{voter}}^2/n_i$ strictly on the $8{,}000$ training profiles yields per-voter rating variance $\hat\sigma_{\text{voter}}^2=\PboVoterSigmaSq$ ($\SD=\PboVoterSD$ points on the $1$--$10$ rating scale). Finite-voter Oracle ceilings rise monotonically from $\rho=\PboOracleTwentyR$ at 20--49 votes to $\rho=\PboOracleFiftyR$ at 50--199, $\rho=\PboOracleTwoHundredR$ at 200+, and $\rho=\PboOracleFiveHundredR$ ($\rho^2=\PboOracleFiveHundredRTwo$, Oracle MAE $\PboOracleFiveHundredMAE$) at 500+ votes (Fig.~\ref{fig:holoface}C). Even where voter sampling error vanishes (500+ votes, median $3{,}526$ raters per portrait), \texttt{HoloFace-Full} leaves $\PboIrreducibleGapFiveHundred$ of true human consensus rank variance unexplained ($\PboIrreducibleGapAll$ across the full test cohort, Oracle $\rho=\PboOracleAllR$). Static photographs and self-reported tags cannot capture dynamic voice, movement, responsiveness, or shared-context conduct ($C_{ic}+H_i$ in Proposition~\ref{prop:twoway}).
\item \textbf{Chemosensory MHC null:} Nor does unmeasured HLA dissimilarity supply a hidden biological shortcut: across $k=63$ human effect sizes synthesized from Havl{\'\i}{\v c}ek et al.~\cite{havlicek2020} and Winternitz et al.~\cite{winternitz2017}, random-effects meta-analytic means include zero across genomic pairing ($z=-0.011$), odor preference ($z=+0.044$), and relationship satisfaction ($z=-0.080$; Fig.~\ref{fig:holoface}D).
\end{enumerate}

\begin{figure*}[!t]
\centering
\includegraphics[width=\textwidth]{performance-fig-face-drivers.pdf}
\caption{\textbf{Canonical 106-landmark Delaunay facial composites and 12-zone Shapley beauty attribution across \PboFaceTotalN{} aligned portraits.} \textbf{(A,~B)}~Super-resolved ($384\times 384$, all cells $n\ge 75$) within-gender bottom-15\%, median (42--58th), and top-15\% composites alongside signed local percentile driver/rejector maps and \PboFaceTrisN-triangle Delaunay log-Jacobian geometry differences for women ($N=\PboFaceEvalNatW$) and men ($N=\PboFaceEvalNatM$). \textbf{(C)}~Exact 12-zone Shapley beauty drivers ($+\Delta$ pct in top 15\%) and rejectors ($-\Delta$ pct in bottom 15\%) with 95\% CIs. \textbf{(D)}~Progressive share of human-jury rank variance ($\rho^2$) from raw UV crop (H1), canonical 106-landmark pure-face model, zero-metadata \texttt{HoloFace-Pure}, and multimodal \texttt{HoloFace-Full} relative to the infinite-voter Oracle ceiling.}\label{fig:face_drivers}
\end{figure*}

\begin{table*}[!t]
\caption{Canonical 12-Zone Facial Beauty Drivers (Top 15\%), Rejectors (Bottom 15\%), and Top-Minus-Bottom Percentile Spread (\PboFaceTotalN{} Portraits).}\label{tab:face_drivers}
\scriptsize
\renewcommand{\arraystretch}{0.97}
\setlength{\tabcolsep}{3.3pt}
\begin{tabularx}{\textwidth}{@{}Lcccccc@{}}
\toprule
\tabinput{performance-face-drivers-table.tex}
\bottomrule
\end{tabularx}
\end{table*}

\begin{figure*}[!t]
\centering
\includegraphics[width=\textwidth]{performance-fig-face-age.pdf}
\caption{\textbf{Super-resolved canonical facial aging composites, anatomy-versus-label decomposition, and 12-zone drift from ages 18--24 to 55+.} \textbf{(A)}~Piecewise-affine Delaunay composites and signed pure-face percentile attribution maps across five adult age brackets (18--24 through 55+, every cell $n\ge 75$). \textbf{(B)}~55+ minus 18--24 rank-shift decomposition separating observed voter drop ($\PboFaceAgeObsYoungW\to\PboFaceAgeObsOldW$th pct in women) into zero-metadata pure facial anatomy ($\PboFaceAgePredYoungW\to\PboFaceAgePredOldW$th pct) plus the non-facial age-label penalty. \textbf{(C)}~Anatomical zone drift between 55+ and 18--24, led by periorbital tear-trough hollowing, nasolabial folding, and jawline contour softening.}\label{fig:face_age}
\end{figure*}

\begin{figure*}[!t]
\centering
\includegraphics[width=\textwidth]{performance-fig-face-world.pdf}
\caption{\textbf{Cross-regional invariance of facial beauty drivers across \PboFaceWorldRegionsN{} world regions.} \textbf{(A)}~Super-resolved regional composites and within-region high-minus-low beauty driver maps across the United States, South America, Mexico/Central America/Caribbean, Canada, Asia-Pacific/Oceania, Europe, and Middle East/Africa (all cells $n\ge 75$). \textbf{(B)}~Pairwise spatial correlation matrix of regional beauty driver maps (lower triangle: women, mean $r=\PboFaceWorldCorrW$, min $\PboFaceWorldMinCorrW$; upper triangle: men, mean $r=\PboFaceWorldCorrM$, min $\PboFaceWorldMinCorrM$). \textbf{(C)}~Observed voter rank versus pure-face predicted rank across regions, separating upload self-selection gaps from invariant anatomical weights.}\label{fig:face_world}
\end{figure*}

\subsection{Canonical 106-Landmark Facial Attribution Across Sex, Age, and World Regions}
To map \emph{where} on the human face visual consensus concentrates without exposing any individual photograph, we registered $\PboFaceTotalN$ front-facing portraits ($N=\PboFaceEvalNatW$ women, $N=\PboFaceEvalNatM$ men; every composite cell $n\ge 75$) onto a canonical 106-landmark, \PboFaceTrisN-triangle piecewise-affine Delaunay atlas~\cite{langlois1990,little2011,deng2019arcface} and computed exact 12-zone Shapley attributions~\cite{lundberg2017shap} from aligned \texttt{DINOv3} + \texttt{ConvNeXt-V2} + \texttt{ArcFace} spatial representations~\cite{woo2023convnextv2,oquab2024dinov2} (live spatial models verified at held-out $\rho=\PboFaceLiveSpearmanW$ in women and $\rho=\PboFaceLiveSpearmanM$ in men; Figs.~\ref{fig:face_drivers}--\ref{fig:face_world}; Table~\ref{tab:face_drivers}).

Three anatomical results emerge. First, beauty drivers and rejectors localize to expressive social-warmth regions---eyes, smile contour, and midface vitality---rather than cranial vault size: in women, \PboFaceTopZoneNameW{} leads total top-minus-bottom spread ($\PboFaceTopZoneSpreadW$ percentile points: $\PboFaceTopZoneDrvW$ driver in top 15\%, $\PboFaceTopZoneRejW$ rejector in bottom 15\%), alongside periorbital eye/iris aperture ($\PboFaceEyesDrvW$ driver), eyebrow arch, vermilion lip and smile contour ($\PboFaceMouthDrvW$), malar cheekbone contour, and tear-trough smoothness; in men, \PboFaceTopZoneNameM{} ($\PboFaceTopZoneSpreadM$ spread: $\PboFaceTopZoneDrvM$ driver, $\PboFaceTopZoneRejM$ rejector), jawline/chin definition, nasal bridge/alar proportion, eyes ($\PboFaceEyesDrvM$), and smile contour ($\PboFaceMouthDrvM$) dominate (Fig.~\ref{fig:face_drivers}A--C; Table~\ref{tab:face_drivers}). Second, across five adult age brackets (18--24 to 55+), pure facial anatomy accounts for gradual periorbital tear-trough and nasolabial softening ($\PboFaceAgePredYoungW\to\PboFaceAgePredOldW$th percentile in women; $\PboFaceAgePredYoungM\to\PboFaceAgePredOldM$th in men), whereas observed voter rankings penalize female profile age more steeply ($\PboFaceAgeObsYoungW\to\PboFaceAgeObsOldW$th; Fig.~\ref{fig:face_age}). Third, across all \PboFaceWorldRegionsN{} world regions, spatial beauty driver maps correlate at mean $r=\PboFaceWorldCorrW$ in women (minimum $\PboFaceWorldMinCorrW$) and $r=\PboFaceWorldCorrM$ in men (minimum $\PboFaceWorldMinCorrM$), even where regional mean voter scores diverge due to upload selectivity (Fig.~\ref{fig:face_world}). Independently, across \PboOkcNoticeN{} self-reported ``First thing people notice about me'' essays in \texttt{OkCupid-JSE}~\cite{kim2015}, smile/laugh (\PboOkcSmilePct{}, $\#1$) and eyes (\PboOkcEyesPct{}, $\#2$) lead every somatic attribute, outranking stature (\PboOkcHeightNoticePct{}) more than twofold and mirroring the 12-zone \texttt{HoloFace-L1} Shapley warmth map.

\begin{figure*}[!t]
\centering
\includegraphics[width=\textwidth]{performance-fig-pov-lifecourse.pdf}
\caption{\textbf{How men and women choose across adulthood, ages 18--75: diverging biological clocks (top row) versus unbroken relational truths (bottom row).} \textbf{(A)}~Realized partner age-gap scissors (partner age minus own age) across $84{,}636$ partnered adults ($37{,}734$ in Kowal et al.~\cite{kowal2025}; $46{,}902$ in Walter et al.~\cite{walter2020}), widening to a $\PboAgeScissorsOld\,\mathrm{yr}$ sex divergence at ages 60--79. \textbf{(B)}~Platform voter rank across \PboHonEligible{} Hot-or-Not profiles versus zero-metadata pure facial anatomy ($5{,}353$ aligned faces), isolating steeper female age penalties beyond anatomical drift. \textbf{(C)}~Explicit ``Want children'' shares among childless users across \PboPovEligible{} Plenty of Fish profiles, crossing at age $\PboPovCrossAge$ (exact age $\PboPovCrossExact$; 25--29 W$-$M gap \PboPovGapYoung{} pp vs.\ 45--49 \PboPovGapOlder{} pp) and replicating across the U.S., U.K., Canada, Brazil, Australia, and open-science San Francisco \texttt{OkCupid-JSE} (age $\PboOkcSfCrossAge\,\mathrm{yr}$~\cite{kim2015}; Table~\ref{tab:pov_geo}). \textbf{(D)}~Constrained 30-point mate-budget trajectories across $67{,}733$ adults~\cite{kowal2025}: kindness ranks $\#1$ in \emph{both} women and men in every age bracket from 18--24 to 60--79, outranking financial resources throughout life. \textbf{(E)}~Serious relationship commitment stays strictly higher in women than men at every adult age ($\PboPovSeriousYoungW\to\PboPovSeriousSenW$ W vs.\ $\PboPovSeriousYoungM\to\PboPovSeriousSenM$ M) alongside rising $>10$-year prior union attainment, proving women's midlife shift reflects childbearing clocks rather than relational exit. \textbf{(F)}~Completed schooling gradient (ages 30--49, POV) showing fertility postponement without loss of childless desire, alongside the partnered happiness lift ($\PboWhyteHappyLiftW$ points W, $\PboWhyteHappyLiftM$ points M on 0--100; $N=\PboWhyteN$~\cite{whyte2021}).}\label{fig:pov}
\end{figure*}

\begin{table*}[!t]
\caption{Cross-National and Open-Science Replication of the Childless Fertility-Intention Crossover Versus Female Relationship Commitment (\PboPovEligible{} POV + \PboOkcN{} \texttt{OkCupid-JSE} Profiles).}\label{tab:pov_geo}
\scriptsize
\setlength{\tabcolsep}{4.5pt}
\begin{tabularx}{\textwidth}{@{}lrrrrrr@{}}
\toprule
\tabinput{performance-pov-geography.tex}
\bottomrule
\end{tabularx}
\end{table*}

\begin{figure*}[!t]
\centering
\includegraphics[width=\textwidth]{performance-fig-gss-fertility.pdf}
\caption{\textbf{Completed-fertility selection floor, \texttt{WORDSUM}/status pairing-versus-births inversion, non-marriage shortfall, and 70-year courtship disintermediation (GSS 1972--2024, $n=\GssN$; Stanford \texttt{HCMST}, $N=\PboHcmstCouplesN$).} \textbf{(A)}~Completed childlessness (ages 45--79) across personal-income fifths: low-income men face a sharp pairing floor ($21.9\%\to 10.5\%$ childless), whereas higher-income women experience elevated childlessness ($9.5\%\to 19.5\%$). \textbf{(B)}~Standardized episode gradients across personal income, occupational prestige, schooling, and objective \texttt{WORDSUM} vocabulary ($g$-proxy): every trait positively predicts male marriage pairing (\GssIncMPair, $\PboGssPrePairM$, $\PboGssSchPairM$, $\PboGssVocPairM\,\SD$) while predicting fewer marital births ($\PboGssVocNumM\,\SD$ on \texttt{WORDSUM}). \textbf{(C)}~Childlessness share of selection opportunity $I$ (rising to $46\%$ in men) alongside the non-marital parenthood gate. \textbf{(D)}~Unchanged stated family-size ideals ($\sim\!2.7$ children) versus the \VgChange{} drop in births per adult aged 25--39 ($\VgMarrShare\%$ in non-marriage). \textbf{(E)}~Partnered probability differences (ages 22--34) for monthly kin contact ($+8.1\,\mathrm{pp}$), weekly worship ($+1.8\,\mathrm{pp}$), neighbors ($-2.6\,\mathrm{pp}$), friends ($-7.5\,\mathrm{pp}$), and social-media hours across 70 countries ($-0.1$ to $-0.6\,\mathrm{pp/h}$). \textbf{(F)}~Stanford \texttt{HCMST} 70-year courtship trajectory~\cite{rosenfeld2019}: mean pre-romance acquaintance collapsed from $\PboHcmstPreEightyYr\,\mathrm{yr}$ (pre-1980) to $\PboHcmstRecentYr\,\mathrm{yr}$ (2012--2017) as online meeting rose to \PboHcmstRecentOnlinePct{} (\PboHcmstRecentAppPct{} dating apps).}\label{fig:gss}
\end{figure*}

\begin{table*}[!t]
\caption{Access--Number Selection Decomposition in Completed U.S.\ Cohorts (GSS 1972--2024, Ages 45--79; Trait $\SD$ per Generation).}\label{tab:gss}
\scriptsize
\setlength{\tabcolsep}{4.5pt}
\begin{tabularx}{\textwidth}{@{}llrCCC@{}}
\toprule
\tabinput{generated/tab_gss.tex}
\bottomrule
\end{tabularx}
\tablenote{Proposition~\ref{prop:episodes} exact partition: $\text{Total}=\text{Access}+\text{Number}$; 95\% bootstrap CIs (2,000 resamples).}
\end{table*}

\subsection{Lifespan Clocks, the Hearth Verb Calendar, Second-Act Re-Entry, and the 94\% Pairing Gate}
Placing $84{,}636$ realized partner age-gap records~\cite{kowal2025,walter2020}, $67{,}733$ mate-budget trajectories~\cite{kowal2025}, \PboHonEligible{} rated portraits, and \PboPovRows{} Plenty of Fish records (\PboPovEligible{} opposite-sex profiles aged 18--79, including \PboVoiceEssaysN{} self-written essays) onto a shared adulthood axis (ages 18--75; Fig.~\ref{fig:pov}; Fig.~\ref{fig:dynamics}D,F; Table~\ref{tab:pov_geo}; Appendix Tables~\ref{tab:pov_resp}--\ref{tab:pov_sens}) separates diverging biological clocks from converging hearth priorities. Across the top row of Fig.~\ref{fig:pov}, biological and casual-market clocks diverge by sex: realized partner age gaps scissor apart to $\PboAgeScissorsOld\,\mathrm{yr}$ by ages 60--79 (Fig.~\ref{fig:pov}A), male-over-female casual sociosexual desire widens from $\PboSoiRatioYoung\times$ at ages 18--24 to $\PboSoiRatioOld\times$ at 60--79 (Fig.~\ref{fig:dynamics}D), crowd photo ratings penalize female age far beyond pure facial anatomy (Fig.~\ref{fig:pov}B), and explicit ``Want children'' shares among childless users flip from women higher at 25--29 (\PboPovWantYoungW{} W vs.\ \PboPovWantYoungM{} M; \PboPovGapYoung{} pp) through an interpolated crossover at age $\PboPovCrossAge$ (exact age $\PboPovCrossExact$) to men higher at 45--49 (\PboPovWantOlderW{} W vs.\ \PboPovWantOlderM{} M; \PboPovGapOlder{} pp; Fig.~\ref{fig:pov}C).

Across the bottom row of Fig.~\ref{fig:pov} and Fig.~\ref{fig:dynamics}F, however, primary relational devotion strengthens across life. First, kindness ranks $\#1$ in \emph{both} sexes in every age bracket from 18--24 to 60--79 (Fig.~\ref{fig:pov}D), and across \PboVoiceEssaysN{} free-text essays, performative ``party/clubbing'' action verbs fall $-\PboVerbPartyDrop$ from ages 18--22 to 65--79 while domestic hearth ``cook/bake'' verbs surge $\PboVerbCookRiseW\times$ in women (\PboVerbCookYoungW{}$\to$\PboVerbCookOldW{}) and $\PboVerbCookRiseM\times$ in men (\PboVerbCookYoungM{}$\to$\PboVerbCookOldM{}), directly echoing the Pleistocene hearth (Fig.~\ref{fig:dynamics}F). Second, women's preference for serious relationship commitment stays strictly above men's in every decade ($\PboPovSeriousYoungW$ vs.\ $\PboPovSeriousYoungM$ at 25--29; $\PboPovSeriousMidW$ vs.\ $\PboPovSeriousMidM$ at 40--49; $\PboPovSeriousSenW$ vs.\ $\PboPovSeriousSenM$ at 60--69; $>10$-year prior union \PboPovTenYrMidW{} vs.\ \PboPovTenYrMidM; Fig.~\ref{fig:pov}E), replicating across all five POV country cohorts and San Francisco \texttt{OkCupid-JSE} (\PboOkcSfGapYoung\,pp at 25--29 flipping to \PboOkcSfGapOlder\,pp at 45--54, crossover age $\PboOkcSfCrossAge\,\mathrm{yr}$~\cite{kim2015}; Table~\ref{tab:pov_geo}), while postgraduate and ``Very Ambitious'' women pair lower current parenthood (\PboPovAmbParentHighW{} vs.\ \PboPovAmbParentLowW) with higher forward child desire (\PboPovAmbWantHighW{} vs.\ \PboPovAmbWantLowW) and partnered adults report $\PboWhyteHappyLiftW$ W / $\PboWhyteHappyLiftM$ M higher happiness ($0$--$100$, $N=\PboWhyteN$~\cite{whyte2021}; Fig.~\ref{fig:pov}F). Third, stratifying midlife re-entrants by marital history confirms Proposition~\ref{prop:lifecycle}(iv): at ages 40--49, divorced women reduce explicit formal marriage intent by $-\PboDivMarryDropW$ relative to never-married single women (\PboDivMarryMidW{} vs.\ \PboSinMarryMidW{}) while raising serious pair-bond commitment demand to \PboDivSeriousMidW{}, whereas widowed men---who lost a cooperative hearth without defection---exceed widowed women in explicit remarriage intent (\PboWidMarryMidM{} vs.\ \PboWidMarryMidW{} at 40--49; \PboWidMarrySenM{} vs.\ \PboWidMarrySenW{} at 60--75; Fig.~\ref{fig:dynamics}F).

Finally, applying Proposition~\ref{prop:episodes} to \GssN{} completed-fertility respondents (ages 45--79) in the General Social Survey (1972--2024~\cite{gss2024}; Fig.~\ref{fig:gss}; Table~\ref{tab:gss}) locates where selection acts today. Within parenthood ($c>0$), family-size differentials (Number) are negative across schooling, objective \texttt{WORDSUM} vocabulary ($g$-proxy, $\PboGssVocNumM\,\SD$ in men), occupational prestige, and income in both sexes. At the marriage pairing gate, however, men face a positive selection gradient across every cognitive and socioeconomic trait---personal income (\GssIncMPair\,$\SD$; Access $\GssIncMAcc\SD$, childlessness falling from $21.9\%$ to $10.5\%$ across income fifths, net $\GssIncMTot\SD$), occupational prestige ($\PboGssPrePairM\,\SD$), schooling ($\PboGssSchPairM\,\SD$), and objective \texttt{WORDSUM} vocabulary ($\PboGssVocPairM\,\SD$, 95\% CI $[+0.001,+0.022]$; Fig.~\ref{fig:gss}A,B)---whereas high-earning women experience higher childlessness ($9.5\%\to 19.5\%$, Access $\GssIncWAcc\SD$; Corollary~\ref{cor:floor}). Across birth cohorts from 1900--1919 to 1960--1979, childlessness rose to account for $46\%$ of men's total opportunity for selection even as non-marital parenthood expanded (Fig.~\ref{fig:gss}C; Appendix Table~\ref{tab:gssgate}). Meanwhile, stated family-size ideals among adults aged 18--44 remained virtually unchanged from 1972 to 2024 ($\sim\!2.7$ children), yet realized births among adults aged 25--39 fell by \VgChange{} children per adult, with $\VgMarrShare\%$ (\VgMarr) accounted for by rising non-marriage (\VgNevOld\%$\to$\VgNevNew\%; Fig.~\ref{fig:gss}D). Consistent with Proposition~\ref{prop:waiting}, seeing relatives monthly predicts $+8.1\,\mathrm{pp}$ higher partnering among young adults (95\% CI $[+2.8,+12.7]$), whereas monthly friend socializing ($-7.5\,\mathrm{pp}$) and daily social-media hours across 70 countries ($-0.1$ to $-0.6\,\mathrm{pp/h}$) do not substitute for kin witnesses (Fig.~\ref{fig:gss}E), paralleling the 70-year Stanford \texttt{HCMST} courtship shift where mean pre-romance acquaintance collapsed from $\PboHcmstPreEightyYr\,\mathrm{yr}$ before 1980 to $\PboHcmstRecentYr\,\mathrm{yr}$ in 2012--2017 as online meeting climbed to \PboHcmstRecentOnlinePct{} (\PboHcmstRecentAppPct{} dating apps; Fig.~\ref{fig:gss}F~\cite{rosenfeld2019}).

\begin{figure*}[!t]
\centering
\includegraphics[width=\textwidth]{performance-fig-confidence-influence.pdf}
\caption{\textbf{Outlier confidence versus gross arrogance: multiplicative prestige, audience priors, the alpha-female double bind, and persuasive coalition calls.} \textbf{(A)}~Across \SdRows{} live speed-dating decisions~\cite{fisman2006}, warm outlier agency ($\ge 8$ partner-rated agency and sincerity) wins \PboWarmOutlierYesW{} (female target) and \PboWarmOutlierYesM{} (male target) ``Yes'' decisions---$\PboWarmHumbleRatioW$--$\PboWarmHumbleRatioM\times$ warm humble low-agency (\PboWarmHumbleYesW{} / \PboWarmHumbleYesM{})---whereas cold high-agency ($\le 6$ sincerity) drops $-\PboColdDropPPW\,\mathrm{pp}$ in women versus $-\PboColdDropPPM\,\mathrm{pp}$ in men ($\PboColdBacklashRatio\times$ steeper female warmth backlash). \textbf{(B)}~Evaluator age and ethnic priors: holding partner-rated traits constant, the penalty on unbacked self-enhancement ($\text{self}-\text{jury}$ agency) surges $\PboOverclaimAgeSurge\times$ from chooser age $\le 24$ ($\PboOverclaimSlopeYoung\,\mathrm{pp/pt}$, null) to $29+$ ($\PboOverclaimSlopeMature\,\mathrm{pp/pt}$) as the raw ambition penalty vanishes ($\PboAmbSlopeYoung\to\PboAmbSlopeMature\,\mathrm{pp/pt}$); across $805{,}460$ prime-age POV profiles, open ``Very Ambitious'' badging spans a $\sim\!2\times$ ethnic gradient (Black \PboEthAmbBlackM{}/\PboEthAmbBlackW{} vs.\ White \PboEthAmbWhiteM{}/\PboEthAmbWhiteW{} and Asian \PboEthAmbAsianM{}/\PboEthAmbAsianW{}), matching live racial choice slopes. \textbf{(C)}~Across 175 countries ($N=\PboKowPrestigeN$~\cite{kowal2025}), warm outlier confidence lifts male partnering $\PboKowWarmOverHumblePP\,\mathrm{pp}$ over humble low-agency and beats cold high-claimers by $\PboKowWarmKindGainM$/$\PboKowWarmKindGainW$ partner kindness and $\PboKowWarmPrqcGainM$/$\PboKowWarmPrqcGainW$ union quality (\texttt{PRQC}), whereas women out-resourcing mates lose $\PboKowOutEarnPrqcDropW$ \texttt{PRQC}. \textbf{(D)}~Asymmetric status threat and executive divorce: outsmarting a partner costs women $\PboOutsmartWomanOnManPP\,\mathrm{pp}$ ($p=0.002$) versus $\PboOutsmartManOnWomanPP\,\mathrm{pp}$ (null) for men, while across \PboExecTotalN{} POV CEOs, founders, and executives vs.\ \PboNonExecPoolN{} non-executives (ages 30--54), female executives show \PboExecDivW{} vs.\ \PboExecDivM{} male executive divorce ($\PboExecDivRatio\times$) and $\PboExecAmbGapPP\,\mathrm{pp}$ lower explicit ambition badging. \textbf{(E)}~Across \PboSccsAllN{} SCCS societies~\cite{kirby2016,boehm1999}, egalitarian bands ($n=\PboSccsEgalN$) enforce \PboSccsEgalCheckedPct{} community checks on leaders, \PboSccsEgalPersuadePct{} persuasion/consultation, and \PboSccsEgalWithdrawPct{} silent walkout removal, whereas class stratification ($n=\PboSccsStratN$) flips authority to \PboSccsStratForcePct{} force. \textbf{(F)}~Proposition~\ref{prop:prestige} closed-form standing curves and $2\times 10^5$-draw Monte Carlo verification (\texttt{E24}, $|\Delta|\le\PboETwentyFourMaxDiff$), alongside the \texttt{OkCupid} Hero's-Journey voice triad ($K=\PboOkcHeroTriadK$ vs.\ $\PboOkcPassiveK$ prompts, $\PboOkcHeroKRatio\times$).}\label{fig:confidence}
\end{figure*}

\begin{table*}[!t]
\caption{Earned Outlier Confidence Versus Gross Arrogance Across Selection and Coalition Arenas (Explicit Evidence Tiers).}\label{tab:influence}
\scriptsize
\renewcommand{\arraystretch}{0.95}
\setlength{\tabcolsep}{3.6pt}
\begin{tabularx}{\textwidth}{@{}>{\raggedright\arraybackslash}p{0.165\textwidth}LLLL@{}}
\toprule
Arena \& Evidence Tier & Earned Outlier Confidence \mbox{($\hat b\uparrow,\delta\approx 0$)} & Gross Arrogance Failure ($\delta>0$ or $\hat p\downarrow$) & Audience Age, Gender \& Ethnic Priors & Alpha Sex Asymmetry \& Co-Protagonist Escape ($h\to 1$) \\
\midrule
\tabinput{performance-influence-regimes-table.tex}
\bottomrule
\end{tabularx}
\tablenote{\emph{Epistemic tiering:} \textbf{Tier~1} statistics are computed deterministically in \texttt{scripts/run.py --study performance-before-outcomes} from primary microdata in this repository; \textbf{Tier~2} citations report external administrative registries, MBA field experiments, and venture-pitch trials~\cite{folke2020,bertrand2015,bursztyn2017,kanze2018,brooks2014,anderson2012,rogers2023}.}
\end{table*}

\section{Outlier Confidence Versus Gross Arrogance: Prestige, Audience Priors, Alpha Asymmetry, and Coalition Influence}\label{sec:prestige}

\subsection{The Multiplicative Prestige Law: Why Warm Outlier Confidence Beats Both Arrogance and Passivity}
Once hominin encephalization unfolded inside coalitionally leveled bands where peers ridicule, desert, or depose coercive bullies~\cite{boehm1993,boehm1999,vonrueden2011}, human status split into two distinct pathways~\cite{henrich2001,cheng2013}: \emph{dominance} (zero-sum intimidation and unbacked self-aggrandizement) and \emph{prestige} (freely conferred deference toward competent contributors). Propositions~\ref{prop:despot}, \ref{prop:delivery}, and~\ref{prop:staging} explain this split through one multiplicative jury equation, extended in Proposition~\ref{prop:prestige} (\texttt{E24}, verified within $|\Delta|\le\PboETwentyFourMaxDiff$ across $2\times 10^5$ Monte Carlo draws; Fig.~\ref{fig:confidence}F): evaluated standing multiplies verified capability/agency $\hat b_g$ by prosocial delivery/warmth $\hat p_g$, $\widehat W_{i,g}=\hat b_g\hat p_g-\gamma_{s,g}\max(b-b_e,0)$. Because $\partial^2\widehat W/\partial\hat b_g\,\partial\hat p_g=1>0$, agency and warmth are strict complements rather than substitutes. Playing small ($\delta=c-b<0$) forfeits entry weight $(1-\rho_b)|\delta|\hat p_g>0$, humble passivity ($b\ll\mu_b$) starves $\hat b_g$, and cold or unbacked boasting ($\delta>0$ or low warmth) collapses $\hat p_g\to 0$.

Tier~1 microdata across both live encounters and global unions confirm this supermodular saddle (Fig.~\ref{fig:confidence}A,C; Table~\ref{tab:influence}). Across \SdRows{} Columbia speed-dating decisions~\cite{fisman2006}, candidates combining outlier partner-rated agency ($\tfrac{1}{2}[\text{intel}+\text{amb}]\ge 8$) with high sincerity ($\ge 8$) achieve ``Yes'' rates of \PboWarmOutlierYesW{} (female targets) and \PboWarmOutlierYesM{} (male targets)---$\PboWarmHumbleRatioW\times$ and $\PboWarmHumbleRatioM\times$ the acceptance of warm, humble low-agency peers (\PboWarmHumbleYesW{} and \PboWarmHumbleYesM{})---whereas equally high-agency candidates rated cold/insincere ($\le 6$) drop to \PboColdArrogantYesW{} and \PboColdArrogantYesM{}. Globally across 175 countries ($N=\PboKowPrestigeN$~\cite{kowal2025}), adults pairing top-tercile capability claims (intelligence, ambition, status $\ge 8/11$) with high warmth ($\ge 9/11$) lift male partnering by $\PboKowWarmOverHumblePP\,\mathrm{pp}$ over warm humble low-agency men (\PboKowWarmPartM{} vs.\ \PboKowHumblePartM{}; vs.\ \PboKowColdPartM{} for cold high-claimers) while securing partners who are $\PboKowWarmKindGainM$ points kinder in men and $\PboKowWarmKindGainW$ points kinder in women ($1$--$11$ scale) and raising realized pair-bond quality (\texttt{PRQC}) by $\PboKowWarmPrqcGainM$ and $\PboKowWarmPrqcGainW$ points above cold high-claimers. Outlier confidence succeeds precisely when high capability is coupled to visible warmth; it turns into gross arrogance the moment $\hat p_g$ detaches from $\hat b_g$.

\subsection{Audience Priors: How Evaluator Age, Gender, Ethnicity, and Culture Shift the Arrogance Cliff}
Why does identical self-promotion succeed in one room yet backfire in another? Proposition~\ref{prop:prestige}(i)--(ii) identifies two audience parameters: evaluator experience/audit depth $m_g$ (which raises overclaiming detection $\rho_\delta(m_g)$) and group humility norm $\kappa_g$. In thin-experience arenas ($m_g\to 1$), observers cannot separate true capability $b$ from bluffing $\delta>0$, temporarily rewarding overconfidence~\cite{anderson2012}; as evaluator experience $m_g$ grows, $\partial^2\E[\widehat W]/\partial\delta\,\partial m_g<0$, turning unbacked boasting steeply negative while preserving returns to genuine ambition.

Live dyadic choices confirm this experience-graded filter directly (Fig.~\ref{fig:confidence}B). Holding partner-rated attractiveness, sincerity, intelligence, fun, and ambition constant, the slope on unbacked self-enhancement ($\text{self agency}-\text{jury agency}$) is statistically null among choosers aged $\le 24$ ($\PboOverclaimSlopeYoung\,\mathrm{pp/pt}$, $t=-0.64$), turns significant at ages $25$--$28$ ($\PboOverclaimSlopeMid\,\mathrm{pp/pt}$, $t=-2.12$), and surges $\PboOverclaimAgeSurge\times$ to $\PboOverclaimSlopeMature\,\mathrm{pp/pt}$ ($t=-4.04, p<10^{-4}$) among choosers aged $29+$, even as the penalty on \emph{actual} partner-rated ambition vanishes from $\PboAmbSlopeYoung\,\mathrm{pp/pt}$ at $\le 24$ to $\PboAmbSlopeMature\,\mathrm{pp/pt}$ ($t=-0.59$) at $29+$. Mature evaluators do not dislike ambitious peers; they punish unverified braggarts. Simultaneously, ethnic and cultural priors shift baseline assertion norms $\kappa_g$: across $805{,}460$ prime-age POV profiles (ages 25--49) spanning nine ethnic cohorts, open ``Very Ambitious'' self-badging spans a nearly twofold cultural range---from Black (\PboEthAmbBlackM{} M / \PboEthAmbBlackW{} W) and Mixed (\PboEthAmbMixedM{} / \PboEthAmbMixedW{}) through Middle Eastern (\PboEthAmbMidEastM{} / \PboEthAmbMidEastW{}), Indian (\PboEthAmbIndianM{} / \PboEthAmbIndianW{}), and Hispanic (\PboEthAmbHispM{} / \PboEthAmbHispW{}) to Caucasian (\PboEthAmbWhiteM{} / \PboEthAmbWhiteW{}) and Asian (\PboEthAmbAsianM{} / \PboEthAmbAsianW{})---aligning with live speed-dating regressions where Black juries positively weight ambition ($\PboRaceBlackAmb\,\mathrm{pp/pt}$) and sincerity ($\PboRaceBlackSinc\,\mathrm{pp/pt}$), Latino juries remain ambition-neutral ($\PboRaceLatinoAmb\,\mathrm{pp/pt}$), White juries penalize both overclaiming ($\PboRaceWhiteOverclaim\,\mathrm{pp/pt}$, $t=-4.57$) and raw ambition ($\PboRaceWhiteAmb\,\mathrm{pp/pt}$), and Asian juries penalize overt ambition ($\PboRaceAsianAmb\,\mathrm{pp/pt}$, $t=-3.18$) while strongly rewarding lively warmth and play ($\PboRaceAsianFun\,\mathrm{pp/pt}$, $t=8.40$).

\subsection{Alpha Male Versus Alpha Female, CEO Culture, and the Universal Hero's-Journey Coalition Call}
Whereas ancestral male pairing selection aligned status directly with hearth provisioning ($\tau_m\approx 0,\gamma_m\approx 0$; Fig.~\ref{fig:gss}B), high-agency and executive women confront a structural double bind governed by prescriptive role-congruity tax $\tau_f>0$~\cite{eagly2002,rudman2001} and insecure-male status comparison $\gamma_f>0$~\cite{fisman2006} (Proposition~\ref{prop:prestige}(iii); Fig.~\ref{fig:confidence}A,D; Table~\ref{tab:influence}). Four Tier~1 results quantify this asymmetry alongside Tier~2 registry and field benchmarks:
\begin{enumerate}[leftmargin=1.4em,itemsep=1.5pt,topsep=2pt]
\item \textbf{Steeper warmth backlash and asymmetric status threat (Tier~1):} Dropping perceived sincerity from $\ge 8$ to $\le 6$ costs high-agency women $-\PboColdDropPPW\,\mathrm{pp}$ in live ``Yes'' decisions versus $-\PboColdDropPPM\,\mathrm{pp}$ for high-agency men ($\PboColdBacklashRatio\times$ steeper penalty). Controlling for all five partner-rated traits, a woman whose intelligence exceeds a man's own self-rating (\texttt{intel $>$ chooser self\_intel}) suffers a $\PboOutsmartWomanOnManPP\,\mathrm{pp}$ ``Yes'' penalty ($t=-3.05$, driven by ambition-insecure men with $\text{self\_amb}\le 7$, where female ambition costs $\PboInsecureMenAmbPP\,\mathrm{pp/pt}$), whereas a man outsmarting a woman incurs zero penalty ($\PboOutsmartManOnWomanPP\,\mathrm{pp}$, $t=+0.24$).
\item \textbf{Executive/CEO divorce, household strain, and strategic self-modulation (Tier~1 + Tier~2):} Across \PboExecTotalN{} CEOs, founders, executives, and business owners versus \PboNonExecPoolN{} non-executives aged 30--54 on POV, female executives ($n=\PboExecWomenN$) exhibit a \PboExecDivW{} divorced share versus \PboExecDivM{} among male executives ($n=\PboExecMenN$; $\PboExecDivGapPP\,\mathrm{pp}$, $\PboExecDivRatio\times$ higher, and $\PboExecMinusNonExecDivW\,\mathrm{pp}$ above non-executive women at \PboNonExecDivW{}), elevated childlessness (\PboExecChildlessW{} vs.\ \PboNonExecChildlessW{}, $\PboExecMinusNonExecChildlessW\,\mathrm{pp}$, opposite in sign to male executives at $\PboExecMinusNonExecChildlessM\,\mathrm{pp}$), and $\PboExecAmbGapPP\,\mathrm{pp}$ lower explicit ``Very Ambitious'' badging (\PboExecVeryAmbW{} vs.\ \PboExecVeryAmbM{}) despite holding identical executive titles; in 175-country union data~\cite{kowal2025}, women out-resourcing mates ($\text{self}>\text{mate}+1$) experience $\PboKowOutEarnPrqcDropW$ lower \texttt{PRQC} ($\PboKowOutEarnCommDropW$ lower commitment) and $\PboKowPhdSingleGapW\,\mathrm{pp}$ higher PhD-versus-high-school singlehood (men $\PboKowPhdSingleGapM\,\mathrm{pp}$). These Tier~1 patterns match Tier~2 Swedish registry causal estimates where promotion to CEO or mayor doubles female divorce without affecting male divorce~\cite{folke2020}, U.S.\ Census household cliffs above a $50\%$ wife income share~\cite{bertrand2015}, and MBA field experiments where single women censor career ambition when observable by male peers~\cite{bursztyn2017}.
\item \textbf{Pleistocene leveling and the Hero's-Journey coalition call across romance, politics, and capital (Tier~1 + Tier~2):} Across \PboSccsAllN{} SCCS societies~\cite{kirby2016,boehm1999,vonrueden2014,vonrueden2016}, egalitarian bands ($n=\PboSccsEgalN$) enforce \PboSccsEgalCheckedPct{} community checks on leaders, \PboSccsEgalPersuadePct{} persuasive or consultative leadership (\PboSccsEgalPurePersuadePct{} pure persuasion vs.\ \PboSccsStratPurePersuadePct{} in class-stratified hierarchies, where authority flips to \PboSccsStratForcePct{} force), $\PboSccsEgalRespectPct$ versus $\PboSccsEgalForcePct$ ($\PboSccsRespectForceRatio\times$) obedience by earned respect over coercion, and \PboSccsEgalWithdrawPct{} removal of arrogant leaders by silent followership walkout (\PboSccsAllWithdrawPct{} across all SCCS; Fig.~\ref{fig:confidence}E). Consequently, legitimate human influence---calling a partner into a life bond, mobilizing a political coalition, or raising entrepreneurial capital---follows the \emph{Hero's-Journey} coalition structure~\cite{campbell1949,rogers2023,willer2009} (Proposition~\ref{prop:prestige}(iv)): demonstrating concrete $K$-domain proof and personal skin in the game while framing the mission as a shared quest ($h\to 1$) that casts partners, voters, or investors as co-protagonists ($\PboETwentyFourHeroRatio\times$ standing gain over self-aggrandizing claims in \texttt{E24}). On \texttt{OkCupid} ($N=\PboOkcN$), writers combining builder-quest purpose (\texttt{essay1}), candid self-admission (\texttt{essay8}), and a co-protagonist adventure call (\texttt{essay9}) without arrogant gatekeeping (\PboOkcHeroTriadPct{}, $n=\PboOkcHeroTriadN$) complete $\PboOkcHeroTriadK/10$ multi-domain prompts versus $\PboOkcPassiveK$ for passive profiles ($\PboOkcHeroKRatio\times$ breadth $K$). Likewise, Proposition~\ref{prop:prestige}(iii) proves that high-agency women cannot escape the double bind by playing small ($\delta<0$, trapped at $0.94$ in \texttt{E24}), but recover near-parity standing ($1.37$ vs.\ $1.40$; Fig.~\ref{fig:confidence}F) through \emph{calibrated competence} ($\delta=0$) paired with \emph{communal/promotion mission framing} ($h\to 1$) and \emph{high-caliber partner or board assortment} ($b_e\ge b_f$)---matching Tier~2 TechCrunch Disrupt venture evidence where investors ask male founders promotion (``how to win'') questions and female founders prevention (``how not to lose'') questions, yet founders who pivot answers to promotion-framed shared vision raise $2\times$--$7\times$ more capital~\cite{kanze2018,brooks2014}.
\end{enumerate}

\section{Discussion, AI Evaluation Parallels, Limitations, and Methods}\label{sec:discussion}

\subsection{Synthesis and Structural Parallel to AI Evaluation}
Across fossil paleoclimate records, small-scale societies, laboratory encounters, and \PboThreePlatformRecords{} digital profiles, a single structural principle recurs: \emph{an evaluation architecture selects for general capacity and primary devotion only when candidates are observed across multiple common domains and unstaged everyday contexts}. When observation collapses to a single curated channel ($k=1,d\to 1$), four distortions follow identically in human dating and leadership markets and in artificial-intelligence benchmark leaderboards~\cite{campbell1979,gao2023,sharma2024,taylor2026iina}:
\begin{enumerate}[leftmargin=1.4em,itemsep=1.5pt,topsep=2pt]
\item \textbf{Single-context saturation ($R^2_{\infty,1}<1$):} Just as thousands of Hot-or-Not voters and our 43-wave \texttt{HoloFace-L1} ensemble leave $\PboIrreducibleGapFiveHundred$--$\PboIrreducibleGapAll$ of true human consensus unexplained because one static portrait omits cross-situational conduct (Proposition~\ref{prop:twoway}), static AI leaderboards saturate on memorized format cues without bounding out-of-distribution generalization.
\item \textbf{Secondary threshold gaming versus breadth dilution:} Just as single-variable height filters generate a $\PboHeightSixZeroRatio\times$ heap at $6'0''$ on POV and a $\PboOkcStrJump$ $5'11''\to 6'0''$ surge exclusively among straight men on \texttt{OkCupid} ($\PboOkcGayJump$ in gay men), optimizing models against narrow scalar reward proxies induces sycophancy and reward overoptimization~\cite{gao2023,sharma2024}; evaluating agents across $K$ diverse, unscripted tasks dilutes domain-specific gaming as $1/K$ (Corollary~\ref{cor:inflation}).
\item \textbf{Capability tests versus unstaged alignment ($d\to 1$):} Announced benchmarks measure maximum capability ($R_b\to 1$) while carrying zero information about unmonitored default rates ($R_p\to 0$; Proposition~\ref{prop:staging}). Bounding deployment reliability requires either unannounced behavioral auditing ($(1-d)n\gg 1$) or architectural execution mediation ($e\to 1$; Proposition~\ref{prop:enforce})~\cite{taylor2026iina}.
\item \textbf{Sycophantic diffidence versus uncalibrated bloviation ($m_g\to 1$):} Under single-turn human preference ratings ($m_g=1$), language models face the same twin failure modes formalized in Proposition~\ref{prop:prestige}---empty overconfident bluffing ($\delta\gg 0$) or servile sycophantic hedging ($\delta\ll 0$)---whereas multi-step execution verification ($m_g\uparrow$) selects for calibrated epistemic confidence paired with cooperative user alignment ($h\to 1$).
\end{enumerate}

\subsection{Scope Limits, Failure Regimes, and Falsifiers}
Every inference in this monograph is subject to explicit boundaries. First, no ancestral heritability, selection gradient, or Paleolithic marriage parameter is estimated; contemporary survey, ethnographic, and platform datasets test only observable comparative statics of the measurement equations. Second, fossil regional series are not ancestor--descendant lineages, Phase III cranial sampling remains geographically uneven, the 930--813\,ka coalescent bottleneck remains actively contested~\cite{hu2023,deng2025,cousins2025}, and coincidence between 1.0--0.78\,Ma cooking evidence and post-0.7\,Ma encephalization is a consistency check rather than causal proof. Third, platform records (HON, POV, \texttt{OkCupid-JSE}) are convenience populations of online users rather than national probability censuses; in the 2021 privacy-revised \texttt{OkCupid-JSE} release~\cite{kim2015}, essay rows are randomly permuted relative to structured demographic rows and ages include bounded privacy jitter, so structured orientation/fertility tables and 10-prompt essay portfolios are analyzed strictly within their respective files without cross-file row linking; cross-sectional age, relationship-stage time-use effort, pre-cohabitation timing, retrospective breakup-initiation reports, executive occupation divorce shares, ethnic self-description rates, and marital-status re-entry profiles combine behavioral responses with cohort and union-survival selection, while causal claims on executive promotion divorce and venture-capital Q\&A framing rest explicitly on cited external Tier~2 registry and experimental studies~\cite{folke2020,bertrand2015,bursztyn2017,kanze2018}. Appendix~\ref{app:falsifiers} (Table~\ref{tab:falsifiers}) pre-commits sixteen discriminating empirical designs and explicit falsifying outcomes across paleoanthropology, juvenile metabolic phylogenetics, experimental partner choice, longitudinal dyadic diaries, behavioral genetics, union-timing survival analysis, and randomized prestige-versus-arrogance pitch framing.

\subsection{Reproducible Pipeline and Privacy Safeguards}
All statistical analyses, Monte Carlo simulations, paleoclimate synchronizations, out-of-fold \texttt{HoloFace-L1} ensemble solves, and vector figures execute deterministically on CPU from SHA-256-verified inputs (\nolinkurl{scripts/run.py} \texttt{--study performance-before-outcomes}, \nolinkurl{scripts/validate.py} \texttt{--source}, and \nolinkurl{scripts/paper.py} \texttt{--study performance-before-outcomes}). The \texttt{LR04} paleoclimate stack (\nolinkurl{evidence/paleoclimate/lr04_benthic_stack.txt}) is public domain; fossil compilations (Gingerich 2022, Will et al. 2021, Du et al. 2018) and modern open microdata (Mehr et al. 2019 \texttt{NHS}, Rosenfeld et al. 2019 \texttt{HCMST}, Kim \& Escobedo-Land 2015 \texttt{OkCupid-JSE}, Whyte et al. 2021, Kowal et al. 2025, Joel et al. 2020, Lidborg et al. 2022, Fisman et al. 2006, \texttt{D-PLACE}/SCCS, GSS) are analyzed from immutable archival releases. On HON, POV, and \texttt{OkCupid-JSE}, minimum cell suppression ($n\ge\PboMinCell$ in demographic and free-text tables; strict $n\ge 75$ individuals per super-resolved canonical Delaunay facial composite) ensures zero individual-portrait or free-text disclosure.
